% !TeX program = pdflatex
% 编译：在仓库根目录运行 .\compile_notebook_tex.ps1 res\G_Framework_Unified_Synthesis_ZH.tex
% 本文件包含全部正文和参考文献，无外部 input/bib 文件依赖。
% 字体沿用《率与商》源文稿：CJKutf8 的 gkai 中文楷体与缩放后的 Bera Mono 等宽字体。
\documentclass[11pt,a4paper]{article}
\usepackage[utf8]{inputenc}
\usepackage[T1]{fontenc}
\usepackage{CJKutf8}
\usepackage[scaled=0.84]{beramono}
\usepackage{indentfirst,titlesec}
\setlength{\parindent}{2em}
\renewcommand{\contentsname}{目录}
\renewcommand{\abstractname}{摘\quad 要}
\renewcommand{\refname}{参考文献与原稿定位}
\renewcommand{\tablename}{表}
\renewcommand{\figurename}{图}
\usepackage[a4paper,left=25mm,right=25mm,top=25mm,bottom=25mm,headheight=15pt]{geometry}
\usepackage{amsmath,amssymb,amsthm,mathtools,bm}
\usepackage{booktabs,longtable,array,tabularx,multirow}
\usepackage{enumitem,setspace,fancyhdr,xcolor}
\usepackage{xurl}
\usepackage[unicode=true,colorlinks=true,linkcolor=black,citecolor=black,urlcolor=black,bookmarksnumbered=true,bookmarksdepth=2,linktoc=all]{hyperref}
\Urlmuskip=0mu plus 1mu
\def\UrlFont{\ttfamily\small}
\DeclareRobustCommand{\gframeworkunderscore}{\raisebox{-0.35ex}{\rule{0.55em}{0.35pt}}}
\makeatletter
\g@addto@macro\UrlSpecials{\do\_{\gframeworkunderscore}}
% 楷体沿用直立字形，西文斜体仍按原环境保留。
\AddToHook{file/c70gkai.fd/after}{%
  \DeclareFontShape{C70}{gkai}{m}{it}{<->CJK*gkaiu}{\CJKnormal}%
  \DeclareFontShape{C70}{gkai}{bx}{it}{<->CJKb*gkaiu}{\CJKbold}%
}
% 顿号后允许断行，使表格中的问题编号和章节编号可在列宽内换行。
\let\gframeworkCJKpunctsymbol\CJKpunctsymbol
\renewcommand{\CJKpunctsymbol}[1]{%
  \gframeworkCJKpunctsymbol{#1}%
  \ifnum#1="01\relax
    \ifnum"\CJK@plane="30\relax\allowbreak\fi
  \fi}
% 仅为中文子字体添加 CJK 映射，保留西文字体自身的 Unicode 映射。
\ifdefined\pdfnobuiltintounicode
  \let\gframeworkCJKaddcmap\CJK@addcmap
  \def\CJK@addcmap#1{%
    \ifx\f@encoding\CJK@enc
      \gframeworkCJKaddcmap{#1}%
      \pdfnobuiltintounicode\font@name
    \fi}
\fi
\makeatother

\hypersetup{pdftitle={G框架的统一数学结构、计算范式与哲学收束},pdfauthor={高政（GaoZheng）},pdfsubject={基于G Framework.txt三十八篇源文稿的中文综合论文},pdfkeywords={G框架,泛逻辑,同伦修复,基础障碍谱系,任务充分商,率与商}}
\setstretch{1.14}
\widowpenalty=10000
\clubpenalty=10000
\displaywidowpenalty=10000
\setlength{\parskip}{0.20em}
\setlength{\emergencystretch}{3em}
\setlength{\tabcolsep}{4.5pt}
\renewcommand{\arraystretch}{1.20}
\setlist{nosep,leftmargin=2em,topsep=0.35em}
\allowdisplaybreaks[2]
\titleformat{\section}{\Large\bfseries}{\thesection}{0.8em}{}
\titleformat{\subsection}{\large\bfseries}{\thesubsection}{0.8em}{}
\titlespacing*{\section}{0pt}{2.0ex plus .5ex}{1.0ex plus .2ex}
\titlespacing*{\subsection}{0pt}{1.4ex plus .3ex}{.8ex plus .2ex}
\theoremstyle{plain}
\newtheorem{theorem}{定理}[section]
\newtheorem{lemma}[theorem]{引理}
\newtheorem{proposition}[theorem]{命题}
\newtheorem{corollary}[theorem]{推论}
\theoremstyle{definition}
\newtheorem{definition}[theorem]{定义}
\newtheorem{assumption}[theorem]{前提}
\newtheorem*{certificate}{证书}
\theoremstyle{remark}
\newtheorem{remark}[theorem]{备注}
\renewcommand{\proofname}{证明}
\newcolumntype{Y}{>{\raggedright\arraybackslash}X}
\newcolumntype{P}[1]{>{\raggedright\arraybackslash}p{#1}}
\pagestyle{fancy}
\fancyhf{}
\fancyhead[L]{\small G框架：统一结构、计算与哲学收束}
\fancyhead[R]{\small 高政}
\fancyfoot[C]{\small\thepage}
\renewcommand{\headrulewidth}{0.3pt}
\renewcommand{\footrulewidth}{0pt}
\title{\bfseries G框架的统一数学结构、计算范式与哲学收束\\[0.65em]\large 从泛逻辑、法空间与同伦修复到任务充分商及率--商统一}
\author{高政（GaoZheng）}
\date{2026年11月1日\quad 综合论述稿}
\begin{document}
\begin{CJK*}{UTF8}{gkai}
\maketitle
\thispagestyle{empty}
\begin{abstract}
本文依据\texttt{G\_Framework.txt}汇集的三十八篇源文稿，对G框架从纯数学基础、应用结构到哲学收束作统一整理。全文以泛逻辑度量与泛迭代为起点，把状态、结构及演化法则共同组织到法空间和GZ-GMS富结构中，以三层联络及主纤维丛广义非交换李代数描述非交换作用、跨层运输和高阶相容性。局部复合产生的Jacobiator-gap、边界缺陷及全局接续障碍，经性质商、语义函子、同伦修复塔与总同调进入统一证书语言；GRL、HACA/PACER、生命总算子丛、物理与工程系统由此取得共同的任务接口。统计格点、LHTS、TPH及轨迹参数同伦把候选生成、延迟信用、残差下降和命中核验接续为可回放的求解流程；PDE替代被表述为任务保持的表达变换与求解迁移，并通过回读、误差和极限条件保持原问题解释。最后，基础障碍谱系把跨学科困难组织成对象、障碍、不变量与见证；《率与商》进一步以积分映射、棱柱同伦、商核下降、执行对应和合法回读，连接传统微积分与语义微积分。本文的哲学归结是：数学的有效性可在规定任务上落实为表示充分、结构相容、路径合法、结果可回读且总成本可核算的联合保持；“率”计算变化与积累，“商”组织任务所需的区分，二者通过任务充分商、微分交换及可验证同伦见证的共同下降形成统一方法。
\end{abstract}
\noindent\textbf{关键词：}G框架；泛逻辑与泛迭代；法空间；PFB-GNLA；同伦修复；总同调；LHTS/TPH；任务充分商；语义微积分；数学有效性
\par\vspace{0.8em}
\noindent\textbf{来源说明：}本文是对所附三十八篇源文稿的综合重写。参考文献编号与附录覆盖索引逐篇对应；日期为整合日期，各源文稿原有日期保留在参考文献中。本文不重复各稿的整段论证，而保存其核心对象、条件、关键推导和理论联系。
\clearpage
\setcounter{tocdepth}{2}
\tableofcontents
\clearpage


\section{研究对象、统一主线与证明约定}
\label{sec:introduction}

\subsection{从多篇形式系统到一个共同问题}
G框架的研究对象同时包括系统的状态、状态之间的合法作用、作用赖以成立的结构，以及改变这些结构的法则。所附文稿从泛逻辑、泛迭代、法空间与主纤维丛出发，逐步发展出语义规范场、同伦修复、统计格点求解和任务商理论；纯数学卷及三部应用卷提供体系骨架，后续形式系统把骨架推进到可构造、可认证和可迁移的接口\cite{S01,S02,S03,S04,S37,S38}。这些对象的共同问题是：当承载、尺度、表示或法则发生改变时，如何仍然保持规定任务所需的含义、可达性、合法性和结果解释？

令$X$表示某一母结构，$\mathcal T$表示规定的任务集合，$q_{\mathcal T}:X\to Q_{\mathcal T}$表示为任务保留必要区分的表示映射，$F$表示演化。全篇围绕以下接口展开：
\begin{equation}
\label{eq:global-spine}
\begin{gathered}
\text{母结构与法空间}
\longrightarrow \text{任务表示与合法作用}
\longrightarrow \text{障碍定位与同伦修复},\\
\text{修复后的候选路径}
\longrightarrow \text{充分商上的求解}
\longrightarrow \text{合法回读与证书核验}.
\end{gathered}
\end{equation}
当$F$能下降到商上时，$q_{\mathcal T}F=\bar Fq_{\mathcal T}$；当严格交换暂不成立而已有同伦见证时，缺陷必须由已定义的边界、微分或高阶操作解释；当所需信息尚不在表示中时，则细化表示并保存来源。这样，基础、方法、应用和哲学具有同一个讨论对象，而无须把每个应用重新写成一套独立理论。

\subsection{概念层次与本文的去重原则}
本文将源文稿按五个相接层次整合：第一层给出泛逻辑、泛迭代、性质组织和法空间；第二层给出联络、非交换算子、障碍及修复；第三层给出语义运行时、开放系统和求解流程；第四层给出物理、生命、认知及工程中的任务实例；第五层给出任务充分商、率--商比较和方法有效性的哲学解释。应用在统一结构中有各自的输入、控制、观测和验收，不以名称相同替代结构映射。

多稿共同使用的同一对象只作一次主定义；后文引用其已建立的结构，补充特定任务的新条件。原稿中“公理式”陈述，能由明确构造推出者写为定理并给出证书与证明；作为结构选择或任务数据者写为定义或前提。证明中的有限迭代采用数学归纳法；涉及无限或超限过程时，另列完备性、连续性、覆盖与极限保持条件。这样的组织使前件可见、依赖可查，并使结论可在新的任务中重新实例化。

\begin{definition}[任务命题与证书的基本类型]
\label{def:task-certificate}
一个任务命题记为$\Gamma\vdash P(x)$，其中$\Gamma$是有来源的前提集合，$x$是带类型的对象，$P$是待验证谓词。其证书记为
\[
\mathcal C=(x,\Gamma,\mathcal R,\pi,\mathcal W,\mathcal B),
\]
其中$\mathcal R$为允许的推演规则，$\pi$为有限依赖图或有限参数化证明，$\mathcal W$为所需映射、同伦、数值界及回读见证，$\mathcal B$为资源和有效域记录。验证器$V(\mathcal C)$逐项核对类型、规则实例、父结论、开放前提及见证。证书中的“已接受”必须指向具体命题和具体有效域。
\end{definition}

\begin{theorem}[证书前缀的条件可靠性]
\label{thm:certificate-soundness}
设一份证书的依赖图有限且无环，每条规则在声明的解释中保持真值，叶节点或者是已核验数据，或者是明确登记的前提。则验证通过的每个结论都在其未解除前提成立时成立；若最终结论所需的开放任务前提为空且计算见证均已完成，则可在声明域内签发该结论。
\end{theorem}
\begin{certificate}[有限推演的归纳证书]
保存拓扑排序$v_1,\ldots,v_N$、各节点的类型与父节点、对应规则及开放前提集合$\mathsf{Open}(v_i)$。数据叶附来源，规则节点附可执行检查，假设解除只能使用已登记的解除规则。所采用的基础解释及规则可靠性是该定理的明示前件。
\end{certificate}
\begin{proof}
对拓扑序前缀长度$n$归纳。$n=1$时，节点为数据叶或前提叶，结论分别由数据核验或所保留前件成立。设前$n$个节点的条件结论成立。第$n+1$个节点的所有父节点均在此前；它所引用的规则保持真值，故父结论在开放前提合取下推出本节点结论。合法的假设解除同步更新$\mathsf{Open}$，未解除的前件不被删除。因此结论及前提登记在每个前缀上共同保持。归纳至$N$得到根结论；若其任务前件均已解决、必需见证均完成，则得到声明域内的实际接受。
\end{proof}

\begin{remark}[数学结构、数值验收与现实命题的接续]
本文的证明给出所声明数学结构和接口上的结论。具体应用通过数据、参数、执行轨迹和外部锚点实例化这些前提；同一证书格式可以承载不同应用，却不代替各应用本身的测量。结构闭合、有限误差接受、概率保证及极限收敛在文中各用相应条件表达。其作用是使“能推出什么”与“需要输入什么”一一对应。
\end{remark}

\subsection{符号统一与来源差异的处理}
全文以$L$及其带下标变体表示相应源稿专用量，$\mathfrak L$或$\mathcal L$表示已注明类型的法空间或法规则族，$\ell_n$表示高阶同伦操作，$d$或$D$表示已声明复形上的微分，$q$表示商或表示映射。D结构的决策组织与总微分$D$按上下文区分；两套$L$编号分别说明为修复层级与语义分辨率，避免将同名层级直接等同。Jacobiator、曲率、任务误差和数值残差各保留自己的类型，只有给出比较映射或不等式后才能互相控制。

综合中采用三项具体校订。其一，完整结合算子代数的交换子满足Jacobi恒等式；非零Jacobiator应在相应投影、截断或高阶括号中解释。其二，基础障碍谱系的希尔伯特问题分类补入源稿另节讨论的第5题，形成1至23的完整索引。其三，若需恢复$p^N$种分型，则二进制信息下界统一写为$\lceil N\log_2p\rceil$，$p=5$仅为特例。另将一般验证、规范化操作与满足$d^2=0$的余链微分分开使用。上述校订用于使各文稿在同一套类型和证明要求下衔接\cite{S01,S27,S37,S38}。


% 本模块覆盖源文献 S01--S16；可拼入使用 amsmath/amssymb/amsthm 的中文论文。
\section{泛逻辑、泛迭代与法空间的统一承载}
\label{fa:foundation}

\subsection{研究对象：对象、规律和规律变换共同进入理论}

G 框架以泛逻辑分析与泛迭代分析为起点：前者规定如何判断可容许性、同一性和路径优劣，后者规定异构对象如何演化以及演化规则如何更新。其统一对象并非某个固定学科的状态向量，而是包含状态、规则、观察、约束及其变换的整体。于是，对象发生变化与对象所依赖的规律发生变化，成为同一体系中不同层次的可描述过程。纯数学卷以法空间、三层联络及 G 代数给出几何与代数承载，后续文献再将这一承载转化为可计算的诊断、修复和执行接口。\cite{S01,S05,S06,S10}

\begin{definition}[异构系统与法对象]
\label{fa:def-lawsystem}
一个异构迭代系统写作
\[
 S_i=(X_i,\mathcal L_i,\Phi_i;\mathcal O_i,\mathcal C_i),
\]
其中 $X_i$ 是状态空间，$\mathcal L_i$ 是局部更新规则，$\Phi_i$ 是由规则生成的演化，$\mathcal O_i$ 是规定观察，$\mathcal C_i$ 是可容许性和保持条件。法对象记录这些数据及其来源；法变换同时给出对象表示、规则和证书的变换。泛逻辑度量比较的是声明任务下的相容性、解释性或路径偏好。源稿的逻辑距离接口取为 $m_{\mathrm{log}}:\mathfrak L\times\mathfrak L\to\mathbb R_{\ge0}$，路径评分取为 $L_D:\Gamma\to\mathbb R$；它们分别记录结构差异与偏好次序，语义均由明确任务和规则固定。
\end{definition}

\begin{theorem}[公理降级：异构系统的共同承载与有限执行保持]
\label{fa:thm-carrier}
给定集合指标的系统族 $\{S_i\}_{i\in I}$，且各系统的有限执行已经良定义。取带标签的不交并
\[
 X=\bigsqcup_{i\in I}(\{i\}\times X_i),\qquad
 \mathcal L=\bigsqcup_{i\in I}(\{i\}\times\mathcal L_i).
\]
将规则 $a\in\mathcal L_i$ 提升为 $\widehat a(i,x)=(i,a(x))$，可构造共同承载，使每个原系统的有限执行、规定观察和合法性判断得到精确保持。
\end{theorem}
\begin{certificate}[带标签嵌入和词长归纳]
保存嵌入 $e_i(x)=(i,x)$、逐生成元交换式 $\widehat a e_i=e_i a$，以及观察和约束的分量定义。归纳谓词为：长度至多 $n$ 的任一合法规则词 $w$ 满足 $\widehat w e_i=e_i w$。
\end{certificate}
\begin{proof}
空词对应恒等映射，故基步成立。若交换式对 $w$ 成立，则
\[
 \widehat{wa}\,e_i=\widehat a\widehat w e_i
 =\widehat a e_iw=e_i a w=e_i(wa),
\]
这里词按先执行左侧原语的约定解释。由数学归纳法，任意有限词均保持。观察和约束在第 $i$ 分量上直接读取原数据，故其值保持。共同承载的构造只证明可统一表达；系统之间的有效转换仍须提供另外的交换见证。
\end{proof}

这一构造说明泛迭代所要求的共同承载可以有明确入口。G 框架的进一步内容，在于共同承载中加入可比较的性质、非交换运输、法曲率和高阶修复，使不同实例之间的关联具有实际数学约束。四维几何与复 Hilbert 空间的耦合，是这一异构表达的典型例子：几何态与算子态通过共同法对象相互影响；其具体物理解释由对应模型及观测拟合承担。\cite{S01,S02}

\subsection{GZ-GMS、法层结构与性质分辨率}

\begin{definition}[广义数学结构与法空间]
\label{fa:def-gms}
GZ-GMS 的统一记账形式为
\[
 \mathfrak G=(\mathsf{Obj},\mathsf{Law},\mathsf{Flow},
 \mathsf{Layer},\mathcal A,\mathcal F_{\mathrm{law}}).
\]
$\mathsf{Obj}$ 承载空间、丛、算子及其范畴；$\mathsf{Law}$ 承载法对象；$\mathsf{Flow}$ 承载可容许流和同伦；$\mathsf{Layer}$ 指定不同分辨率可见的性质；$\mathcal A$ 是层间联络；$\mathcal F_{\mathrm{law}}$ 是法曲率。GZ-LS（高氏法空间）可写为 $(\mathfrak L;J,\mathrm{Loc},\Sigma)$，其中 $J$ 记录度量，$\mathrm{Loc}$ 记录局部化规则，$\Sigma$ 记录触发细化或结构转换的阈值。GZ-LT、GZ-LOC、GZ-TLC 分别指法变换、法联络及三层联络接口；GZ-LCurv 对应法曲率族 $\mathcal F_{\mathrm{law}}$。
\end{definition}

这里的富结构性是可核查的数据承载：将同一对象的算子作用、组合顺序、观察投影、性质档案与证明见证共同保存。固定其中某些字段、忽略某些方向或按任务关系取商，便得到严格代数、传统几何或离散程序等具体表示。来源不同的构造可以共享同一承载，但是否可互换仍由它们实际保持的字段决定。\cite{S01,S09,S10}

\begin{proposition}[性质塔的规范投影]
\label{fa:prop-propertytower}
设性质抽取器 $P_k:X\to V_k$ 满足 $P_k=r_kP_{k+1}$。定义 $x\sim_k y\iff P_k(x)=P_k(y)$，则存在唯一满射
\[
 q_{k+1,k}:X/{\sim_{k+1}}\longrightarrow X/{\sim_k},\quad
 [x]_{k+1}\longmapsto[x]_k,
\]
且这些映射按层复合。对于由性质塔保持的同伦见证，具有该见证的对象在每层上同类；若性质塔另有完备的见证重建程序，反向也成立。
\end{proposition}
\begin{certificate}[因子化与重建接口]
逐层提供 $r_k$、性质的同伦不变性证明；若要从全部性质相同恢复同伦，另提供重建程序和正确性证书。
\end{certificate}
\begin{proof}
$P_{k+1}(x)=P_{k+1}(y)$ 推出 $P_k(x)=r_kP_{k+1}(x)=r_kP_{k+1}(y)=P_k(y)$，故类映射良定义且由代表元唯一确定。逐层投影都把 $[x]$ 送到更粗的 $[x]$，所以复合一致。见证保持各 $P_k$ 给出正向结论；反向由附加的重建证书给出。
\end{proof}

性质塔说明“同一性”可以随任务分辨率组织：粗层合并暂时不可区分的对象，细层恢复任务需要的差异。文献中的 CH 分层和连续统崩溃—重建（GZ-CBT），在本文以这种语义分辨率变化、法对象转换及证书重建来统一表述。经典集合论的基数层和高阶结构的性质层分别保留各自的判定口径；经典 CH 及其否定的相对一致性仍按原形式系统陈述，语义分层则记录增加了哪些观察、规则和见证。相同同调群只是指定不变量相同；需要完整同伦同一性时，必须保存相应见证或完备性条件。\cite{S01,S06,S10}

\subsection{生成式集合观、单子接口与证书化同一性}

单子 $(T,\eta,\mu)$ 用“进入上下文—组合上下文—消去嵌套”组织对象生成。它使性质、来源和构造证据能够随对象一起生成，契合 G 框架将规律构造过程纳入对象的要求。性质变化由带高阶分量的态射表达，结构同一性由可组合的见证表达；严格同一性是其中见证可退化为恒等的特例。\cite{S09,S10}

\begin{proposition}[公理降级：有限词生成器的单子律]
\label{fa:prop-monad}
令 $T(X)=X^*$ 为有限词集合，$\eta_X(x)=[x]$，$\mu_X$ 为有限词列表的串接。则
\[
 \mu_XT(\eta_X)=\mathrm{id}_{T(X)},\quad
 \mu_X\eta_{T(X)}=\mathrm{id}_{T(X)},\quad
 \mu_XT(\mu_X)=\mu_X\mu_{T(X)}.
\]
故有限算子词提供生成式集合观的一个明确单子实例。
\end{proposition}
\begin{certificate}[嵌套词与次序保持]
保存每个词的有限原语序列、空词及串接算法；所有压平操作均保持原语先后次序。
\end{certificate}
\begin{proof}
两个单位律分别把词中每个元素包装后拆开，或把整个词包装后拆开，所得仍是原词。对三层词，等式两侧都按原来的外层、内层次序枚举同一原语序列。以外层词长归纳，空词时相同，追加一项只在两侧追加相同序列，故结合律成立。
\end{proof}

有限词单子只是最小可执行实例。超限生成、完备滤过及带高阶相容关系的生成器，还须另外给出极限对象与运算良定义性。由此，“公理降级”的实际含义是：把能够由明确构造推出的结构律改写为定理，并显式保留这些定理所依赖的基础规则和输入条件。

\section{三层联络、G 代数与决策—修复结构}
\label{fa:algebra}

\subsection{几何、参数与法规则的联络}

在局部可微模型中，三层基底取为 $X_{\mathrm{tot}}=M\times\Lambda\times\mathfrak L_0$。主丛 $\widehat P\to X_{\mathrm{tot}}$ 上的总联络 $\mathcal A$ 同时描述几何方向、参数方向和法规则方向的运输，曲率
\[
 \mathcal F=\mathrm d\mathcal A+\tfrac12[\mathcal A,\mathcal A]
\]
则包含各层内部及层间混合方向的非交换信息。固定参数与法对象得到传统几何切面；允许参数变化得到连接族；将参数进一步提升为算子泛函，便进入可变泛函—算子版本。只有总连接与各投影相容时，层间运输才可交换。\cite{S01,S05}

以空间坐标 $x$ 与法参数 $w$ 为例，可记
\[
 A=A^{(x)}(w)+A^{(w)},\qquad
 F=F_{xx}+F_{xw}+F_{ww}.
\]
其中 $F_{xw}$ 登记“先改变状态再改变规则”与反向操作的差异。主丛的群作用用于组织可逆的表示冗余；执行动作一般形成幺半群，允许不可逆作用。二者在同一总结构中相容，但不因此被视为同一个代数对象。

\begin{theorem}[严格算子恒等式与曲率的分别承载]
\label{fa:thm-strict}
在定义域对复合封闭的结合算子代数中，分次交换子严格满足分次 Jacobi 恒等式。对于连接 $\nabla$，
\[
 [\nabla_X,\nabla_Y]-\nabla_{[X,Y]}=R(X,Y),\qquad
 \mathrm d_A F=0.
\]
非零曲率与严格 Jacobi 恒等式可以同时成立。
\end{theorem}
\begin{certificate}[共同定义域、交换子展开与连接数据]
登记算子共同作用域、结合复合、分次符号、连接及曲率的定义。投影或截断后的新括号须另行登记。
\end{certificate}
\begin{proof}
将分次交换子代入三重 Jacobi 和，六种三算子复合按分次符号两两抵消。第二式中的 $R$ 按连接交换子的定义得到。把 $F=\mathrm dA+A\wedge A$ 代入 $\mathrm d_AF=\mathrm dF+A\wedge F-F\wedge A$，利用 $\mathrm d^2=0$ 与外微分乘法律逐项抵消，即得 Bianchi 恒等式。因此曲率记录平行运输的非平直性，而完整算子交换子的 Jacobi 和仍为零。
\end{proof}

\begin{remark}[版本整合口径]
原始材料使用“Jacobiator-gap”统摄多种闭合缺口。本文统一规定：严格交换子使用上述恒等式；非零 Jacobiator 则属于投影、截断、有效括号或一般高阶结构，并由其具体定义计算。曲率、三元残差和上同调障碍之间的联系，均通过所声明的映射和见证建立。
\end{remark}

\subsection{四种 G 代数及高阶一致性}

PFB-GNLA 即“基于主纤维丛的广义非交换李代数”。纯数学卷给出四种相互衔接的表述：严格版本保留精确算子关系；同伦版本以 $\ell_3,\ell_4,\ldots$ 吸收低阶失配；可变泛函—算子版本把法规则本身作为可变化的算子对象；元数学原版则从泛逻辑与泛迭代重建整个法层结构。它们在 GZ-GMS 中共享母结构，在各自满足投影和提升条件的区域内相互对位。\cite{S01,S07,S09}

采用悬移余代数上的次数一余导子 $Q=Q_1+Q_2+\cdots$，高阶一致性统一写为
\[
 Q^2=0.
\]
该式在原次数下给出 $\ell_1^2=0$、$\ell_1$ 对 $\ell_2$ 的相容性，以及 $\ell_3$ 对 $\ell_2$ 的 Jacobi 失配的修复。固定符号约定后，第 $n$ 阶义务写成
\[
 \delta \ell_n=-\mathcal O_n,\qquad \delta\mathcal O_n=0,
\]
其中 $\delta$ 是低阶微分诱导的算子微分，$\mathcal O_n$ 由已有低阶项计算。对固定复形，存在 $\ell_n$ 当且仅当 $[\mathcal O_n]=0$。因此，缺口的数值大小用于选择修复顺序，而障碍类是否为零决定当前结构内的修复可解性。\cite{S07,S08,S09}

\begin{proposition}[新增边缘算子的修复判据]
\label{fa:prop-edge}
设 $\delta^2=0$，$\delta\mathcal O=0$。在当前复形中，方程 $\delta h=-\mathcal O$ 可解，当且仅当 $[\mathcal O]=0$。若扩展到 $(C',\delta')$，则需重新验证 $\delta'^2=0$、原任务保持及 $[i\mathcal O]=0$，才能在扩展中修复该障碍。
\end{proposition}
\begin{certificate}[障碍代表元与填补见证]
保存 $\mathcal O$ 的具体余链值及闭性计算；成功时提供 $h$ 并逐项检查 $\delta h+\mathcal O=0$。扩展时另保存 $i:C\to C'$ 的链映射性质、保护字段和合法回读。
\end{certificate}
\begin{proof}
若 $h$ 存在，则 $\mathcal O=-\delta h\in\operatorname{im}\delta$，故其上同调类为零；反之类为零即意味着存在该原像。对扩展复形使用同一论证即可。扩展是否保持旧任务不由上同调等式单独决定，因此需同时核验所列保护字段。
\end{proof}

“边缘算子”在此表示补入结构接续缺口的算子。变分构造提供其一个来源：在具有合适分次 Poisson 括号的 BV 模型中，$Q=\{S,\cdot\}$，且 $\{S,S\}=0$ 推出 $Q^2=0$。这将“修复算子从何而来”落实到作用量、主方程和计算域。Bianchi、Jacobi 与高阶一致性在这一意义下形成几何、代数和证书的对应接口。\cite{S08,S09}

GZ-OHU 负责将局部修复组织为算子—同伦升级塔。其完整证书应同时包含后继阶段的合法附加、已有任务保持、待修复义务的覆盖调度及极限阶段的相容性。有限工程截断保存残余缺口和精度预算；超限完成保存相应序数阶段的映射和极限见证。结构存在、普适因子化、数值最优和有限时间求解各有独立前提，本文在后续修复章节分别落实。

\subsection{D 结构的决策语义与修复深度}

D 结构把环境、历史和权重基准组织成路径选择系统，DERI 在该接口中承担由环境与历史拟合权重基准的任务。记
\[
 w=\mathsf{Fit}_D(E,H),\qquad
 L_D(\gamma;w)=\sum_{j=1}^{m}w_jd_j(\gamma).
\]
评分诱导路径偏序，执行后的新经验又更新 $(E,H,w)$。D-Decision 指给定基准上的路径选择，D-Depth 指基准和结构的递归更新深度。源文献将其与 Jacobiator-gap 的缺口修复对应，关键共同对象是“有定义域、有偏好、有更新规则的路径语义”，而非评分数值本身。\cite{S05,S06}

\begin{theorem}[D 评分与缺口评分的路径商对应]
\label{fa:thm-decision}
在同一路径集 $\Gamma$ 上，设 $f$ 与 $g$ 诱导完全相同的严格排序。则在实际取值域之间存在唯一严格递增双射 $\varphi:f(\Gamma)\to g(\Gamma)$，满足 $g=\varphi f$。因此，在已有法层实现见证的区域内，D 评分与证书评分可以在路径偏好商上相互转换。
\end{theorem}
\begin{certificate}[同序性、同值性与法层实现]
验证任意两路径的严格比较一致；若把 $g$ 声称为几何缺口的函数，另外提供该缺口与路径的实现映射。
\end{certificate}
\begin{proof}
定义 $\varphi(f(\gamma))=g(\gamma)$。若 $f(\gamma)=f(\eta)$ 而 $g(\gamma)\ne g(\eta)$，则 $g$ 的某个严格比较经同序性推出 $f$ 的严格比较，矛盾。故映射良定义。严格同序性给出严格递增及单射，定义给出满射，且其取值已由 $g$ 唯一确定。
\end{proof}

对于更新塔 $D_{n+1}=\mathsf{Upd}_D(D_n,\gamma_n)$，严格下降的自然数排名通过定理~\ref{ss:termination}给出有限终止证书。若使用一般良序排名，则得到无无限严格下降链的结论；若讨论超限构造的稳定序数，则还须给出极限阶段规则。修复广度记录参与作用的缺口分量或新增方向，修复深度记录接续层级，二者共同描述求解结构的组织成本。

\section{统一应用模板：路径积分、物理模型与生命总算子丛}
\label{fa:applications}

\subsection{非交换算子词与 GRL 路径测度}

三个应用卷共享如下模板：先指定法对象与任务观察，再给出可组合算子、代价和证书，最后在合法执行路径上构造测度。离散执行采用算子幺半群；KAT（含测试的 Kleene 代数）进一步将动作、条件测试、选择与迭代纳入同一程序语言。GRL 将策略理解为这些法空间路径上的分布，既可更新数值参数，也可选择操作序列和规则组合。\cite{S02,S03,S04}

设候选路径有限，先验测度 $p_0(\gamma)>0$，$\Gamma_{\mathrm{ok}}$ 为已经通过硬约束与证书的非空集合。取 $\beta>0$，定义
\[
 Z_\beta=\sum_{\gamma\in\Gamma_{\mathrm{ok}}}
 p_0(\gamma)e^{-\beta S(\gamma)},\qquad
 q_\beta(\gamma)=\frac{p_0(\gamma)e^{-\beta S(\gamma)}}{Z_\beta},
\]
其中 $S$ 可以合并任务成本、残差、不确定性和资源使用。硬证书以可行域裁剪表达；软偏好通过有限代价表达，二者的职责不同。

\begin{proposition}[GRL 有限路径积分的变分证书]
\label{fa:prop-gibbs}
对于支撑在 $\Gamma_{\mathrm{ok}}$ 上的概率分布 $q$，令
\[
 \mathcal F_\beta(q)=\sum_\gamma q(\gamma)S(\gamma)
 +\beta^{-1}\sum_\gamma q(\gamma)
 \log\frac{q(\gamma)}{p_0(\gamma)}.
\]
则 $q_\beta$ 唯一最小化该泛函，且
\[
 \mathcal F_\beta(q)=-\beta^{-1}\log Z_\beta
 +\beta^{-1}D_{\mathrm{KL}}(q\|q_\beta).
\]
\end{proposition}
\begin{certificate}[有限路径、代价和归一化]
保存合法路径集合、每条路径的 $S$、$p_0$ 及 $Z_\beta>0$；核验所有质量都分配给已认证路径。
\end{certificate}
\begin{proof}
把 $\log q_\beta=\log p_0-\beta S-\log Z_\beta$ 代入相对熵并求和即得恒等式。相对熵非负，等号仅在 $q=q_\beta$ 时成立，故得唯一最小值。
\end{proof}

这一有限版本明确连接路径优化、统计力学和证书化策略搜索。连续路径积分另需测度或离散极限；量子路径中的复振幅、干涉与酉演化按其量子模型处理，不能仅由概率 Gibbs 权重替代。GRL 的通用程序表达由有效编码支持，具体加速及所需硬件由编码、访问、采样和验证成本共同决定。\cite{S02,S16}

\subsection{物理、量子计算与复杂材料的共同接口}

物理应用以主丛连接承载规范场，以法参数承载有效模型及其变化，以路径作用量组织演化和控制。其四条应用线索为：一般相对论与量子理论的法层耦合；GRL 的策略及随机控制；量子信道、测量投影、电路路径和经典—量子混合求解；超导、超流及二者协同的多尺度材料构造。它们共同使用“状态—法规则—操作路径—证书”的接口，具体守恒律、作用量、测量结果和实验可辨识性各自实例化。\cite{S02}

对于复杂材料，材料成分只是总状态的一部分。微结构、外场、边界、几何、制造、控制和扰动恢复共同决定可行域；总 Jacobiator-gap 聚合指定层次上的闭合残差，实验观测承担模型校准。由此，逆向设计转为在明确目标和可复现条件下搜索可执行构造路径。室温超导、超导—超流协同及相关计算载体在源稿中提供研究任务与候选模型；每个实际实例均以其物理评价器和验证记录签发。

\subsection{LBOPB：七算子幺半群、六观察框架与多尺度生命对象}

生命总算子主纤维丛（LBOPB）写为
\[
 \mathcal B_{\mathrm{life}}=
 (M_{\mathrm{life}},P_{\mathrm{life}},G_{\mathrm{life}},
 \omega_{\mathrm{life}};\{\mathcal M^D\}_{D\in\mathcal D}).
\]
基底记录声明分辨率上的生命状态，纤维记录内部配置，法联络组织跨视角运输，七个算子幺半群提供可组合的干预及演化原语。\cite{S03}

\begin{center}
\begin{tabular}{ll}
\hline
记号 & 在统一生命模型中的职责\\
\hline
PEM & 病理演化与疾病进展\\
PRM & 生理调节与稳态维持\\
TEM & 毒理效应与风险累积\\
PKTM & 药代输运、分布及清除\\
PGOM & 药物基因组与相关程序化作用\\
PDEM & 药效响应和效应组合\\
IEM & 免疫激活、抑制及相互作用\\
\hline
\end{tabular}
\end{center}

七个算子域与六个观察框架承担不同职责。六个框架分别观察病理、药效、生理调节、毒理、药代以及合并的药物基因组—免疫信息。对同一总轨迹 $\gamma$，它们给出 $\pi^r\gamma$；源稿所谓全息对应，指多视角由共同轨迹诱导，并在充分性成立时共同决定该轨迹的规定任务信息。同一疾病也可以分别以生理、病理或干预参数为基底，其余自由度并入纤维；不同丛表示应通过观察和运输交换式对齐。

\begin{proposition}[多框架一致性及其可恢复范围]
\label{fa:prop-multiview}
设总演化为 $F_a:X\to X$，第 $r$ 个框架的演化为 $F_a^r$，且 $\pi^rF_a=F_a^r\pi^r$。则任意有限操作词在所有框架上保持一致。若联合观察 $\Pi=(\pi^r)_r$ 在规定任务商上单射，则各框架联合结果能够唯一恢复该任务商中的总状态。
\end{proposition}
\begin{certificate}[逐原语交换与联合观察充分性]
保存全部 $\pi^r$、生成元交换式和目标任务；恢复时提供 $\Pi$ 在任务商上的单射或明确解码器。
\end{certificate}
\begin{proof}
以词长归纳，恒等词显然保持；添加原语 $a$ 后，交换式把 $\pi^rF_aF_w$ 转为 $F_a^rF_w^r\pi^r$，故有限词保持。若两总状态给出同一联合观察，商上单射性使它们落在同一任务类，故所述任务信息唯一恢复。
\end{proof}

微观分子对接及分子动力学、介观组织与分室模型、宏观病理和药效算子，通过投影、提升和法联络组成多尺度模型。逆向药物构造据目标效应选取算子词，再连接分子生成与模拟后端；虚拟人体将这些过程组织为随时间更新的大系统，支持模型内虚拟试验和长期演化研究。七域属于可扩展的工程基准，其充分性总是相对于目标、分辨率、可观测性和干预域判定。知识压缩则把已验证的跨域对应编译为可复用算子和证书，避免每次从全部微观细节重新推导。\cite{S03,S13}

\section{HACA、PACER 与语义规范场的可执行闭环}
\label{fa:semantics}

\subsection{从字符、词包到性质与决策}

HACA 是分层代数认知架构，组织符号、语义和决策三层；PACER 是词包对齐的压缩—展开推理器，执行压缩、对齐、展开、再压缩的循环。其多尺度操作域写作
\[
 \Omega=\Omega^{(0)}\sqcup\Omega^{(1)}\sqcup
 \Omega^{(2)}\sqcup\Omega^{(3)},
\]
依次对应字符、词包、句式及段落路径。性质层从这些表达抽取声明任务的语义承诺，其意向性幂集只选取可观察、同伦稳定并可在法空间实现的性质子族；同伦层保存对改写、组合和扰动的稳定见证，跨域函子层组织类比和隐喻。语义函子代数作为 PFB-GNLA 的受约束表示，使“含义怎样组合”和“组合缺口怎样修复”进入同一结构。\cite{S04,S11}

两个耦合的 GRL 路径积分分别处理词法行动成本和段落语义可容许性：表达流畅或词法代价较小，不能单独签发语义接受；语义合格也不自动确定最有效的词包组织。合法段落需同时满足两类约束及其层间交换。MDQ 即微分动力量子，在语义实例中是可独立检查、组合和回放的最小语义步。原稿中的注意力、时间和频率可以由路径占用测度与资源预算具体化。

\subsection{部分规范化、最小标准族与歧义收缩}

\begin{definition}[语义规范化接口]
\label{fa:def-normalizer}
设表达域为 $\mathsf{WP}$，规范形域为 $\mathsf{NF}$。一个部分规范化接口为
\[
 F:\mathsf{WP}\to\mathsf{NF}\sqcup\{\bot\},\qquad
 F(x)=\begin{cases}\mathrm{nf}(x),&\chi(x)=1,\\\bot,&\chi(x)=0.\end{cases}
\]
$\chi$ 是可核验的接受谓词；$\mathrm{nf}$ 给出其接受域中的规范形；$\bot$ 表示当前接口不能签发。若称其为范畴意义的函子，还须给出对象、态射及恒等和复合的保持数据。
\end{definition}

语义规范场可由一族标准化接口组织。所谓标准，是规则、版本和验证器可追溯；所谓最小，是在指定分离或生成任务下不可再删除，而非脱离任务的绝对最小。有限候选族可以逐项删去冗余接口而得到一个不可再约族；该族一般未必唯一，也未必具有最小基数。表达的语义签名与规范形根据被保留的任务性质决定，风格和表面措辞可以在同一签名内变化。\cite{S11,S12}

\begin{theorem}[公理降级：可靠过滤的有限候选收缩]
\label{fa:thm-filter}
设初始候选集 $C_0$ 有限，真实可接受类 $s_*\in C_0$。每轮使用可靠测试 $V_t$，满足 $V_t(s_*)=1$，并令
\[
 C_{t+1}=\{s\in C_t:V_t(s)=1\}.
\]
则 $s_*\in C_{t+1}\subseteq C_t$；严格收缩至多发生 $|C_0|-1$ 次。若测试族能区分 $s_*$ 与所有其他候选，且区分测试最终均被执行，则最终候选为 $\{s_*\}$。
\end{theorem}
\begin{certificate}[保留真值、有限集与公平测试调度]
保存初始覆盖、每轮测试和被删除候选的失败见证；唯一性要求补充测试分离能力及调度覆盖证书。
\end{certificate}
\begin{proof}
基步由 $s_*\in C_0$ 成立。若 $s_*\in C_t$，可靠性给出 $V_t(s_*)=1$，故 $s_*\in C_{t+1}$，并由定义有包含关系。严格收缩每次至少删去一个非 $s_*$ 的元素，故次数有界；最终所有其他候选均遭其区分测试删除时，只剩 $s_*$。
\end{proof}

该定理给出重复扫描的实际价值：新测试必须减少尚存歧义，或明确返回当前分辨率下的未决集合。对于未覆盖输入、证书冲突和有限预算耗尽，运行时保留原因与后续义务，不把未决项重新命名为已验证结论。

\subsection{压缩—展开闭环、双重指纹及多实现一致性}

\begin{theorem}[语义保持的 PACER 回环]
\label{fa:thm-pacer}
设 $\pi:X\to S$ 为语义投影，$G:S\times\Sigma\to X$ 为带风格参数的展开器，并在合法域内满足
\[
 \pi(G(s,\sigma))=s.
\]
则从任意 $x_0$ 出发，递推 $x_{n+1}=G(\pi(x_n),\sigma_n)$ 时，有 $\pi(x_n)=\pi(x_0)$。若第 $n$ 轮引入合法语义变换 $T_n$，则相应投影严格依次为 $T_{n-1}\cdots T_0\pi(x_0)$。
\end{theorem}
\begin{certificate}[投影回读与每轮变换]
逐次保存输入签名、展开参数、输出回读和所使用的合法语义变换；任何回读不一致都终止该轮接受。
\end{certificate}
\begin{proof}
无变换时，$\pi(x_{n+1})=\pi G(\pi(x_n),\sigma_n)=\pi(x_n)$，由归纳法得到不变量。存在 $T_n$ 时，同式变成 $\pi(x_{n+1})=T_n\pi(x_n)$，再归纳展开即可。
\end{proof}

MDQ 最小运行时记录作用片段、类型、验证谓词、规范化重写和证书。有限格式可采用正则表达式或自动机作为检查见证；更丰富的语义由相应判定程序及其证明接口承担。上层提案器输出 MDQ 调用序列 $\mathrm{Comp}_m(x)$，统一执行器给出 $\mathrm{Eval}(\mathrm{Comp}_m(x))$。多实现可移植性的条件是：各实现被接受时，都回到同一规定语义 $\pi(x)$。语义指纹记录规范化内容，证书指纹记录证据、验证器与策略版本；哈希用于身份和版本关联，内容正确性仍由实际验证器承担。\cite{S12}

HACA/PACER 的完备性以三层分别表达：LA 要求至少一条到达终态的合法路径；LB 进一步要求全部合格路径共享规定的性质核心；LC 再要求声明的规范域内路径具有共同的同伦正规型。三级条件用于表达单路径存在、共同语义与规范路径稳定性，须分别提供见证。阅读路径证书、占用测度、缺口残差及版本差异共同支持长文档比较、审计与回滚。\cite{S04}

由此，白盒语义系统的闭环是：提出候选、构造规范形、检查证书、执行路径、回读语义、记录残差并修复结构。语言模型可以作为提案器、路由器、片段标注器、策略器或风格展开器；评价其作用时，使用误收、漏收、延迟、成本及回退次数。是否接受则由声明任务和证书接口决定。这使模型改进与语义标准可以各自更新，同时保留跨版本比较的共同基准。

\section{开放工程、两层增强与 ORL 的边界}
\label{fa:openengineering}

\subsection{从生命总算子丛迁移到物理与工程总算子丛}

PTOPB 与 ETOPB 分别强调物理总算子丛与工程总算子丛。它们沿用 LBOPB 的总对象—算子词—路径积分—证书模板，同时把材料、结构、制造、控制、验证和运维纳入联合状态。令
\[
 S=X_{\mathrm{phys}}\times Z_{\mathrm{eng}},\qquad
 s_{t+1}=\pi_t\delta_t(s_t),
\]
其中 $\delta_t$ 是外部来源的扰动，$\pi_t$ 是内部可选构造。物理/工程属于作用侧分类，外生/内生属于来源及可控性分类，两种分类交叉使用：物理过程也可以被控制，工程链路也可能承受外部扰动。\cite{S13,S14,S15}

\begin{proposition}[生成元兼容推出整条设计路径兼容]
\label{fa:prop-compile}
给定状态适配器 $h:X\to Y$ 和原语编译 $c:\mathcal M_X\to\mathcal M_Y$。若每个可执行原语 $a$ 满足 $h(a x)=c(a)h(x)$，且 $c$ 保持词连接，则对每个有限设计词 $w$ 有
\[
 h(w x)=c(w)h(x).
\]
若 $c$ 在所比较的有限路径集间为双射、逐路径先验权重相同，且两侧代价满足 $|S_Y(c(w))-S_X(w)|\le\varepsilon$，则对应路径积分的自由能差不超过 $\varepsilon$。
\end{proposition}
\begin{certificate}[适配器、原语表与统一误差界]
保存 $h$、原语编译表、逐原语交换计算、合法域与代价误差；自由能比较要求对应路径测度及支撑同步。
\end{certificate}
\begin{proof}
空词时交换式成立；追加一个原语后，先用原语交换式，再用归纳假设，得词级交换。代价界给出
$e^{-\beta\varepsilon}Z_X\le Z_Y\le e^{\beta\varepsilon}Z_X$，取 $-\beta^{-1}\log$ 得自由能界。
\end{proof}

该命题是虚拟环境、策略、证书和多保真评价器复用的共同依据。实际迁移先检查跨域适配是否保持任务，再复用搜索；同一数学模板不会自动赋予不同物理后端相同的有效域。

\subsection{可达性诊断、控制共设计与非交换治理}

逆向设计先比较路线的可达下界，再在可行路线中优化。记路线 $k$ 的允许路径为 $\Gamma_k$，诊断值为
\[
 D(k)=\inf_{\gamma\in\Gamma_k}S(\gamma).
\]
扩大合法候选族不会抬高该下确界，但缩小模型误差和取得可达见证仍需实际工作。目标与代价的对应还需区分“存在零代价实现”“下确界为零”和“已有正下界证书”：前者提供达到目标的路径，后两者分别记录逼近能力和不可逾越的残差。将控制变量与材料、几何共同设计，可避免因过早冻结控制而产生伪不可达判断；将残差、不确定性和复现证书纳入评价，可发现路线代价下降但实际解释力未改善的偏离。\cite{S13}

\begin{proposition}[相同材料投影允许不同治理结果]
\label{fa:prop-material}
存在保持材料变量 $m$ 的两个算子 $a,b$，其执行顺序改变任务结果。故在含此类实例的任务族中，完整可行性不能仅由材料投影决定。
\end{proposition}
\begin{certificate}[两个构造算子及可复算反例]
令状态为 $(m,u,v)$，取 $a(m,u,v)=(m,1,v)$，$b(m,u,v)=(m,u,u)$，目标为 $v=1$，初态为 $(m,0,0)$。
\end{certificate}
\begin{proof}
按先 $a$ 后 $b$ 执行得 $(m,1,1)$，按先 $b$ 后 $a$ 执行得 $(m,1,0)$；二者材料投影都为 $m$，目标真值不同。若目标完全因子化为材料谓词，则相同 $m$ 必给出相同真值，与计算矛盾。
\end{proof}

室温超导路线在源稿中又按目标规格分为常压稳态室温工程化、中高压高临界温度、常压中高临界温度或亚稳态维持三类。低阻旁路、屏蔽和主动控制属于相应工程效用目标；严格超导目标则使用其物性与稳定性谓词签发。目标层次在总模板内分别比较，可避免把不同的成功条件合并成同一分数。

这一非交换反例保留了工程材料稿件的重要结论：构造顺序、控制与保障可以成为决定性任务变量。对于工业超导、芯片、发动机及基础设施，具体是否存在这样的决定性路径，由各自的物理和工程算子实例化；“治理”在数学上指组织这些可控操作及其证书，使系统在规定扰动中维持或返回任务域。

\subsection{相内优化、包络迁移与三类变换}

设工程模型形成族 $\mathcal B_\lambda$。相内优化在固定 $\lambda$ 下改变代表元和控制参数；包络迁移改变允许的自由度、算子域、约束或修复复形。若某个缺口分量对全部允许原语都不变，则对任何有限算子词仍不变，因而非零初值给出可计算的不可作用下界。若在已知可行零点附近，连续控制到缺口向量的导数满秩，则隐函数或次浸没接口提供局部修正方向。固定复形中的非零障碍类不能靠更换同类代表元消除，因而必要时必须进入第二层搜索。对子空间、纤维包络和丛塔的选择，是将哪些变量提升为基底、哪些变量保留为纤维的显式组织方式。\cite{S13,S14}

与此对应，需区分三类变换：规范变换改变同一结构的表示；同一规范结构中的非平凡演化改变状态或连接；扇区切换改变有效结构与算法适用域。扇区切换需要同步搬运状态、保护量、证书和算法条件。总守恒量约束整个母结构，局部子系统仍可通过边界交换而改变。工程生命周期由这种带条件的有向转换图表示，逆向设计则在前向可达集与目标的反向可达集之间寻找交会。

OpenRA 等实时对抗沙盘在源稿中承担低成本可执行实例：观测不完全、资源竞争、非交换动作、环境变化与多时标控制都可显式记录。子丛解耦先固定接口轨迹，再并行求解局部任务；拼接需要接口与资源约束共同满足。相同组织模板再连接地磁扰动下的电网节点、量子计算节点抗退相干、工业设计及自动化科研。统一的是任务和修复接口，具体安全域和真实效果由领域模型承担。\cite{S14,S16}

\subsection{开放完备性、未知事件与语言扩展}

开放系统除状态变化外，还可能增加此前未被表示的观察、约束或算子。记 $\mathcal E$ 为允许扰动类，$\mathcal A=\{s:J(s)\preceq J_{\max},\ \mathrm{Cert}(s)=1\}$ 为任务接受域。工程完备性要求：对 $\mathcal E$ 中每次已声明扰动，存在合法修复词使系统在规定长度、资源、副作用和时限内回到 $\mathcal A$。其等级可记为 $(\kappa,\ell)$，分别描述扰动开放程度和所需修复层级；该等级是指定模型上的描述参数。\cite{S15}

固定观察的信息边界由命题~\ref{pc:information-limit}应用于交互历史接口得到：两个世界若给出相同初始记忆、相同已读取观察及相同随机输入，算法就给出相同执行。若两世界所需的成功动作集合不相交，该接口不能同时保证两者成功。其具体证书是一对同历史、异目标的世界及其合法动作记录。

在可测口径下，这对应固定事件族 $\Sigma_0$ 对更新和复合封闭；要依赖原本不可测的新事件，必须引入新观察或细化表示。不断缩短回合、增加内部搜索或多代理互审，都无法单独生成两个世界之间尚未进入输入的区分信息。外部真值通道、独立验证器和受保护的规则扩展因此是开放科研的实际组成部分。多代理一致性可以提高候选质量，证书接受仍须接到同一明确任务的证据基础。\cite{S14,S15,S16}

\subsection{幂域纤维、离散规范场与 ORL}

为表达宏态下的微观搜索，取 $\widetilde\Theta=\Theta\times G^k$，投影 $p(\theta,z)=\theta$。映射
\[
 \theta\longmapsto\{\theta\}\times G^k
 \in\mathcal P^+_{\mathrm{fin}}(\widetilde\Theta)
\]
在 $G$ 有限非空时为单射；每个宏态提升为带可枚举微观自由度的纤维。源稿中的十维格点是 $k=10$ 的实例。递归纤维扩展、调度层数与有效一步执行可以分别登记：高阶搜索组织不等于每次物理执行必须采用高阶原语。\cite{S16}

在完备度量承载（源稿称卡丘空间）上，选定分辨率与网格可以组织离散骨架 $\Gamma=(V,E)$。可逆传输由边连接 $U(e)$ 表示，非可逆执行则继续保留在适当的范畴或幺半群中。正实权重场可写为 $A(e)=-\beta\log w(e)$、$U(e)=e^{A(e)}$；带方向的回路周期满足
\[
 \log\mathrm{Hol}(c)=\sum_{e\in c}\epsilon_eA(e).
\]
固定生成树后，基本回路给出独立周期测试；指定弦边的修复可以消去这些周期。对一般结构，还需核验修复是否保护原任务。统计权重调整改变路径测度；删边或把概率置零则改变可达支撑，二者在求解意义上不同。

ORL 在这些材料中指面向开放系统、具有观察—行动—残差—结构修复闭环的学习和控制。它以可回放路径为运行对象，以证书为接受接口，以新增算子和性质承载开放扩展。相对于给定观察接口 $\mathcal O$、验证器 $\mathcal V$、算子覆盖和预算，源稿将其与可审计白盒 AI 及 $\mathrm{AGI}^{\mathrm{cert}}_{\mathcal V,\mathcal O}$ 对位；等价成立于明确接口和完备性条件内。

有限路径采样可以用当前轨迹与必要统计量替代整条路径空间的显式枚举，形成时间换空间的计算方式；其收益由混合、命中、验证及输出成本核算。微观扭曲频率增加的是单位物理时间内执行的步数，温度和能量则决定稳态偏好。独立证据通道与同源失效隔离还可进入工程保护设计，例如源稿将光域通信和存储作为对特定电磁耦合失效的分离方案；其有效性由实际耦合域、辅助设备和能量供给条件确定。由此，理论统一、运行求解与工程保护在同一证书链上对接。\cite{S16}


\section{联络传输、语义编码与可认证运行时}
\label{ss:semantic}

在统一法空间中，几何、语义与程序执行通过共同的可达对象接续。这里需要同时保存对象如何被编码、路径如何被选择以及结论如何被验证。文献\cite{S19,S20,S22}分别承担传输编码、语义运行和扩展签发三项接口；其共同内容可以压缩为\emph{同一任务对象在多种表示之间的保真传递}。

\subsection{可达终点、路径支撑类与同伦正规形}

\begin{definition}[联络可见的路径商]
固定主丛 $P\to M$、联络 $A$、起点 $p_0$ 及可容许基底路径集 $\Gamma_A$。假定每条路径都存在唯一水平提升，记提升终点为 $e_A(\gamma)$，并定义
\[
 \gamma\sim_A\eta\quad\Longleftrightarrow\quad e_A(\gamma)=e_A(\eta),
 \qquad Q_A:=\Gamma_A/\sim_A.
\]
可达终点集为 $F_A=e_A(\Gamma_A)$。路径积分的支撑信道以 $Q_A$ 为指标；同伦扭曲正规形集 $T_A$ 则由一个已给定的规范化规则选取每个等价类的唯一代表。支撑信道记录传输对象，数值权重另由作用量与测度给出。
\end{definition}

\begin{theorem}[三重编码的同一对象]
\label{ss:three-encoding}
在上述条件下，存在显式双射
\[
 F_A\cong Q_A\cong T_A.
\]
若另给三侧的拓扑结构，且编码、解码连续，则双射提升为同胚；若只提供互逆至同伦的映射及见证，则得到相应同伦等价。
\end{theorem}
\begin{certificate}
保存水平提升规则、终点等价判据、正规化算法 $N_A$ 及 $N_A^2=N_A$。拓扑结论还保存连续性或两侧同伦见证。曲率非零时，以端点固定的普通路径同伦替换 $\sim_A$，须另验联络运输在该同伦下保持。
\end{certificate}
\begin{proof}
映射 $\bar e_A:Q_A\to F_A$ 由 $[\gamma]\mapsto e_A(\gamma)$ 给出。等价关系的定义保证其良定义与单射，$F_A$ 的定义保证满射。映射 $[\gamma]\mapsto N_A(\gamma)$ 由正规化对等价类的恒定性良定义；每类唯一代表给出逆映射。后两种加强分别由所提供的连续逆或同伦逆直接成立。
\end{proof}

这一结果使联络提升、离散路径积分和同伦扭曲共享一个对象入口。对有限合法路径族，沿用命题~\ref{fa:prop-gibbs}的路径权重，在支撑类上比较作用量，再回落为有来源的终点代表。若最小作用量类唯一，低温权重集中于该类；若并列，则保留最优类集合或使用明确的并列规则。每一层终点还可封装为下一层输入，使三重编码沿修复塔递推；这种递推依赖层间投影与编码交换，不能由单层双射替代。

\subsection{规范语义、逻辑占位与文本表达}

文献\cite{S20}把语义系统组织为三层：规范语义塔保存定义、量词、推理与证书；逻辑占位塔保存待表达结构；文本层产生可读表述。设 $S_b$ 是上下文 $b$ 上的规范语义对象，$P_b$ 是占位对象，$Y_b$ 是表面文本，则基本接口写为
\[
 S_b\xrightarrow{\iota_b}P_b\xrightarrow{R_b}Y_b
 \xrightarrow{\operatorname{parse}_b}P_b.
\]
占位对象包括量词范围、谓词依赖、对象类型、证书引用、失败分支、来源接口和表达参数。问答内核据此被定义为上下文修复塔中的合法路径类，回答是该类经占位层的表达投影。领域知识则以认证纤维、局部约束、专用算子和验证器进入插件；主基底保留任务定义与最终验收规则。

\begin{proposition}[被接受表达的语义保持]
\label{ss:render}
设 $\iota_b$ 在规定语义域内可逆，解释映射为 $I_b:S_b\to\mathcal M_b$。验证器仅当
\[
 \operatorname{parse}_b(y)=\iota_b(s),\qquad
 \operatorname{Verify}_b(s,c)=1
\]
时接受文本 $y$。若验证器对证书 $c$ 的判定可靠，则每个被接受文本表达 $I_b(s)$，且其被签发结论具有相同证书。
\end{proposition}
\begin{certificate}
记录规范对象标识、占位编码、解析结果、证书依赖及验证器版本。渲染器的随机性、语言风格和候选数都列入运行记录；接受规则本身由确定规则执行。
\end{certificate}
\begin{proof}
由解析等式及 $\iota_b$ 的可逆性，解码结果唯一为 $s$；再经 $I_b$ 得到同一意义对象。验证器可靠性把证书的结论传给 $s$。两项联合保证被接受表达同时保持意义与证明来源；候选文本如何产生不影响这一推导。
\end{proof}

原文所称“白盒主权”据此具有具体含义：求解对象、可用规则、选择依据、约束检查与结果签发均由可检查接口控制；大语言模型可以作为可替换的表达或候选组件。一个插件若只增加新类型及新规则，并证明原有对象上的旧判断不变，则是保守扩张。每次持久化保存规范对象、路径类、上下文、版本和证书，而后续调用先核验上下文兼容性，再重用已认证结构。由此，知识积累成为结构与证书的可追踪积累。

主基底与扇区插件共用六层签名：算子与自由度、纤维化状态空间、允许路径类、作用量选类、正规形收束、证书构造与规则签发。可调用规则由末端的正规形和证明产生，搜索权重则承担候选选择。这个顺序使新规则的取得具有可回溯的求解和验证来源。

\subsection{真理锚点、拒绝签发与修复入口}

文献\cite{S30,S31}进一步规定任务锚点
\[
 a=(\text{对象域},\text{形式规则},\text{评价标准},\text{证据接口},\text{验证器}).
\]
事实记录是证据输入，形式可靠性与适用域决定何种结论可由这些输入推出。运行时若未经声明更换标准，输出可能在新标准下局部自洽，却不再回答原任务。这里“锚点保真”要求每次标准扩展都有明确映射，且在旧任务上的解释保持；它同时允许新增证据与合法细化。

\begin{theorem}[有限语义链的锚点保持：公理的定理化]
\label{ss:anchor}
令初始状态 $s_0$ 满足锚点谓词 $A(s_0)$。若每个被接受修复步 $r_n$ 都附有
\[
 A(s_n)\wedge V(r_n,s_n)=1\ \Longrightarrow\ A(r_n(s_n)),
\]
且 $s_{n+1}=r_n(s_n)$，则所有有限前缀保持 $A$。
\end{theorem}
\begin{certificate}
归纳证书为初值见证、各步验证记录、适用域和锚点映射。未经证明的额外标准须作为新任务字段登记，不能暗中成为旧结论的新前件。
\end{certificate}
\begin{proof}
基步由 $A(s_0)$ 成立。若 $A(s_n)$ 成立，则本步验证记录与蕴含式给出 $A(s_{n+1})$。归纳得到任意有限前缀的保持性；证明只使用实际给定的保持见证，未把保持要求本身当作已完成结果。
\end{proof}

扩展签发协议把命题分为来源直接证明、补充证书后证明和未证明。直接证明对应新增证书包为空的特例，因而包含于扩展签发域。对未闭合陈述，系统输出缺口
\[
 r=(\text{缺失定义},\text{类型失配},\text{断裂推理},
       \text{证据不足},\text{资源边界}),
\]
再选择补定义、补对照、补算子、补验证器或补适用域。所谓概念降级，是将原命题改写为证书实际支持的条件化命题，同时记录两者关系；语义附加则为目标对象增加必要语境，若采用伴随语言，还须给出相应自然对应。拒绝在此成为修复入口，修复无从完成时签发当前边界。

这一规则有两个互补方向：不把缺少证明的强结论当成可调用事实，也不把已在既定标准下闭合的结论无故削弱。证书检查应定位具体断点，并给出有内容的修复或边界。正常放行要求残差已按声明阈值处理；保留残差的分支必须显式标记，避免语义孔洞被文本流畅性掩盖。

标准本身也纳入审计。令 $\mathcal S^*=\mathcal S\cup\operatorname{Audit}(\mathcal S)$，其中审计项覆盖标准选择、阈值设定、可比对象和标准切换。对于已观测的特征交集 $E=\bigcap_iE_i$，联合概率应按条件概率链
\[
 q_j=\mathbb P_j(E_1)\prod_{i\ge2}\mathbb P_j(E_i\mid E_1\cap\cdots\cap E_{i-1})
\]
计算。若证书给出 $N_jq_j<1$，得到该标准及样本尺度下期望出现次数小于一的稀有事件判定。原文称之为标准簇反身闭合的低基率交集点；它不需要把相关特征误当独立，也不允许在结论形成后无记录地改换阈值。这个例子将强结论是否成立还原为标准、概率和来源的具体计算。

\subsection{修复塔的可恢复表示}

文献\cite{S22}的单一实数编码可以作如下统一表达。若每层数据已被忠实编码为实数 $x_n$，先用 $\nu(x)=\frac12+\pi^{-1}\arctan x$ 将其归一化，并取规范二进制位 $b_{n,m}$。对自然数配对双射 $\langle n,m\rangle$，定义
\[
 E((x_n))=\sum_{n,m\ge0}2b_{n,m}\,3^{-(\langle n,m\rangle+1)}.
\]
三进制位只用 $0,2$，其恢复规则逐位还原 $b_{n,m}$，再逆归一化恢复各层。若两输入有不同的首个编码位，该位的差不能被其后所有位抵消为同一允许编码，故 $E$ 单射。编码说明可数塔存在统一的数学载体；读取无限精度、验证全部层级和输出完整结构的费用仍由实际访问接口承担。这使“结构可以统一表示”与“结构可以廉价求解”各自保留精确位置。

集合元素的结构表示还有一个直接模型：对任意集合 $S$ 取离散空间 $S_{\mathrm{disc}}$，则 $H_0(S_{\mathrm{disc}};\mathbb Z)\cong\mathbb Z^{(S)}$，每个 $s\in S$ 对应一个独立生成元 $[s]$。其证书是零阶链群、边界为零及基元识别；不同元素落到不同基元，因此表示忠实。这给出“一切皆结构”的弱数学版本：对象可作为结构中的生成位置被追踪。实数总编码仅进一步适用于已经满足逐层可数编码接口的塔，两种结论各自保留基数及访问条件。

\section{障碍层级、离散修复与宏观--微观总同调}
\label{ss:repair}

\subsection{可修复障碍、广义分形与超限单子}

文献\cite{S18}将可修复障碍作为对象总域，将具有递归尺度结构的广义分形放在其中的特定子域。记 $\mathcal R$ 为带修复见证的障碍对象类，$\mathcal F\subseteq\mathcal R$ 为进一步带兼容尺度递归的对象类，则分形性描述生成形态，修复性描述能否把指定障碍化为边界。有限一次填补即可消去的对象可属于 $\mathcal R$ 而不满足所要求的非平凡尺度递归，因此二者不应合并命名；传统自相似构造在满足修复条件时作为 $\mathcal F$ 的特例。

该区分与“超限单子”即高氏单子的定位相接。其基元携带可执行相互作用、非交换算子、性质投影与修复自由度；“卡丘胚”即非交换投影核主丛胚，则组织这些投影及其核。以离散统为起点，修复生成新的合法接续，再经几何实现得到连续或时空表象。原文把这条路径写为离散统、卡丘胚、修复塔、几何实现的生成链；其中每一箭头都须给出承载结构与投影证书。静态冻结可以恢复旧结构的核，动态演化则继续保留被旧投影隐藏的修复方向。

轨道商与核商分别承担两类归并：前者归并指定作用下等价的表示，后者消去指定映射不可见的方向。若交换子落入投影核，非交换结构在该投影下呈现交换化；该性质由 $q([a,b])=0$ 检查。莱布尼茨“无窗”单子的对照在此只涉及所设定的因果接口：无交互的离散范畴可嵌入更丰富的相互作用范畴，而含锚映射、联络和修复链的基元拥有额外可达结构。这样，基础改写的内容落在对象及态射数据上。

\subsection{同调填补与离散谱证书}

\begin{theorem}[一次离散修复的精确效果]
\label{ss:cell-repair}
设有限维实内积链复形在各阶正交归一基下的矩阵为 $B_2:C_2\to C_1$、$B_1:C_1\to C_0$，满足 $B_1B_2=0$。给定闭链列 $Z$，即 $B_1Z=0$，置
\[
 B_2^+=[B_2\ Z],\quad B_1^+=B_1,
 \quad \Delta_1=B_1^{\mathsf T}B_1+B_2B_2^{\mathsf T}.
\]
则修复后仍为链复形，并且
\[
 b_1^+=b_1-\bigl(\operatorname{rank}B_2^+-\operatorname{rank}B_2\bigr),
 \qquad \Delta_1^+=\Delta_1+ZZ^{\mathsf T}.
\]
全部一阶同调被消去时，$\Delta_1^+$ 正定；固定 $C_1$ 内积时，其特征值不小于原特征值。
\end{theorem}
\begin{certificate}
保存 $B_1,B_2,Z$、矩阵秩、同调基、内积和谱计算。几何附胞腔还需验证各列闭链可被声明的附着映射实现；任务观察的保持另验相应交换式。来源为\cite{S21}。
\end{certificate}
\begin{proof}
由 $B_1[B_2\ Z]=[0\ 0]$ 得到链复形合法性。再用
\[
 b_1=\dim\ker B_1-\operatorname{rank}B_2
\]
相减即得同调维数公式。直接相乘得拉普拉斯增量。对任意 $v$，
\[
 \langle v,\Delta_1^+v\rangle
 =\|B_1v\|^2+\|(B_2^+)^{\mathsf T}v\|^2.
\]
其核为闭链中与边界正交的谐和代表，故核维数等于 $b_1^+$。当该维数为零时正定；半正定增量与最小最大刻画给出特征值单调性。
\end{proof}

同调、扇区数量与谱隙构成三种互补检查。高维复内积空间与算子代数可在此基础上承载最小量子数据 $(\mathcal H,\mathfrak A,U(t))$；指定自伴生成元 $H$ 后，$U(t)=e^{-itH}$ 给出演化。源稿的量子--几何接口同时保留几何重建、维数、度量与观测的附加条件，使代数可运行性和物理解释各自具有明确入口。

\subsection{Jacobi--Bianchi 障碍与有限终止}

文献\cite{S23,S24}把局部截面不能延拓、Jacobi 相容失败和离散连接的曲率接续失败转化为可登记障碍。在具有相应余链模型的任务中，三阶缺陷 $o_3$ 满足 $do_3=0$，其障碍类为 $[o_3]\in H^3(C)$；若存在 $h_2$ 使 $dh_2=o_3$，修复便以该边界见证消去此类。具体的 Jacobiator、Bianchi 缺陷与上同调类之间须由映射定义相接，不能只凭名称对齐。

\begin{theorem}[有限缺陷下降给出终止]
\label{ss:termination}
设每个已认证状态 $X_n$ 有整数缺陷量 $\mu_n\ge0$。只要 $\mu_n>0$，系统均能在给定预算内产生合法修复 $R_n$，且
\[
 \mu(R_nX_n)\le\mu_n-1.
\]
则至多 $\mu_0$ 个有效修复步到达 $\mu_N=0$；若 $\mu=0$ 的定义逐项覆盖全部登记缺陷，则终端代表满足相应闭合条件。
\end{theorem}
\begin{certificate}
证书包括初始缺陷表、每步修复、任务保持、严格下降和终端复核。这里“零障碍”指该表规定的障碍；新增任务引入的新缺陷需进入新的表。
\end{certificate}
\begin{proof}
对 $n$ 归纳有 $\mu_n\le\mu_0-n$，直到终止。若经过 $\mu_0$ 步仍未终止，下一步要求负整数缺陷，与非负性矛盾。终端逐项复核把计数归零传递为登记谓词的合取，故所声明闭合成立。
\end{proof}

有限修复骨架把超限原型投影到当前可访问的有效层。预算增加时，若可行集嵌套 $\mathcal P_B\subseteq\mathcal P_{B'}$，则
\[
 \inf_{\gamma\in\mathcal P_{B'}}\mathcal A(\gamma)
 \le\inf_{\gamma\in\mathcal P_B}\mathcal A(\gamma).
\]
这说明更大预算扩大可比较路径；对当前集合的全局最优，须与更深修复空间中的最优分别标记。完整终端证书还记录所选路径、提升障碍及费用，从而同时回答终点为何可用与路径如何到达。

\subsection{连续基底与离散修复塔}

为避免不同论文的层号混用，本节将文献\cite{S26}的层级记为 $\mathsf L_1,\mathsf L_2,\mathsf L_{\ge3}$：$\mathsf L_1$ 是可数离散索引，$\mathsf L_2$ 是连续几何基底，$\mathsf L_{\ge3}$ 是覆盖在该基底上的离散修复纤维塔。双离散对仍可数，但可数索引的无限兼容序列可编码连续对象。具体地，给定非空闭集
\[
 K_{\sigma|n+1}\subseteq K_{\sigma|n},
 \qquad \operatorname{diam}K_{\sigma|n}\longrightarrow0
\]
及完备底空间，嵌套收缩得到唯一极限点；若每个点均有兼容提升分支，则编码覆盖整个基底。连续几何、离散证书与高阶修复因此可以共同存在。

这些修复层实现了第~\ref{fa:foundation}节的 CH 语义分层：连续基底承载基数及几何表示，高阶纤维承载性质、路径和障碍修复。

路径层先剔除违反类型、约束和证书的分支，再在余下空间中定义作用量。原文的“视域加权测地线”就是这一允许路径空间内按观测权重选出的极小代表；存在性由适当的紧性与下半连续性保证，唯一性则由额外的严格凸性或正规化选择保证。删除非法路径和改变路径权重分别登记，避免在评价变化中遗漏实际可达性的改变。

层数增大与障碍量下降是两条独立进程。若 $r_n\le C\rho^n$、$0<\rho<1$，只能先推出 $r_n\to0$；要获得极限截面，还需层间相容、紧性或完备性，以及缺陷泛函在所用拓扑下的闭性。若同时优化连续参数，则修复误差和优化误差可写成
\[
 E_{n,k}\le E_{\mathrm{repair}}(n)+E_{\mathrm{opt}}(n,k),
\]
分别核验两项，再选对角序列 $k=k(n)$ 使两项归零。这正是存在性修复与可行域最优化的双重收敛接口。

\subsection{宏观--微观总微分与时间结构}

设宏观空间为 $M$，微观纤维为 $F$。文献\cite{S27}采用
\[
 C_{p,q}=C_p(M)\otimes C_q(F),\qquad
 D_0(x\otimes y)=\partial_Mx\otimes y+(-1)^p x\otimes\partial_Fy.
\]
联络的水平--垂直分解给出两个方向，跨层耦合体现在分量的相互依赖；总复形进一步保存耦合对闭合性的影响。取形式修复塔
\[
 D^{[n]}=D_0+\sum_{r=1}^{n}\varepsilon^rK_r.
\]

\begin{theorem}[总同调的逐阶闭合：公理的归纳证明]
\label{ss:total}
若 $D_0^2=0$ 且每阶满足
\[
 D_0K_n+K_nD_0+\sum_{i=1}^{n-1}K_iK_{n-i}=0,
\]
则 $(D^{[n]})^2=0\pmod{\varepsilon^{n+1}}$。在完备分离的形式幂级数承载上，$D=D_0+\sum_{r\ge1}\varepsilon^rK_r$ 满足 $D^2=0$，故总同调 $H(C,D)$ 良定义。
\end{theorem}
\begin{certificate}
保存每个 $K_r$ 的次数、适用域、逐阶系数方程和滤过。若将形式参数取实数值，还需算子级数收敛及连续乘法证书；该分析实现与形式闭合分别登记。
\end{certificate}
\begin{proof}
基步为 $D_0^2=0$。从 $n-1$ 到 $n$，原来低于 $n$ 次的系数不变，第 $n$ 次系数恰为所列方程左端，因此前 $n$ 阶继续为零。归纳得到任意有限截断。形式完备承载中每一系数只涉及有限项，且全部系数为零，故 $D^2=0$；于是 $\operatorname{im}D\subseteq\ker D$，同调商有定义。
\end{proof}

每步方程是否有解由相应算子复形的障碍类决定。若不存在所需 $K_n$，该类作为剩余障碍返回；若引入新纤维或新算子，则同时核验旧任务观察的保持。宏观离散分型与微观连续跃迁幅由总复形连接，单一扇区标签不足以替代完整微观态。

时间在该体系有两种角色：修复链的累积代价 $T_n=\sum_{j<n}c_j$ 给出离散过程时间，几何实现上的时间函数给出因果曲线的有向顺序。单独的复内积空间没有指定动力学生成元；加入 $H$ 后才得到上述 $U(t)$。BA 交替链保存微观承载到宏观实现、再反馈到下一轮承载的双向作用。时间属性可先赋予 $\mathbb Z_2$ 分次，但只有奇微分平方为零后才产生相应分次同调。

源稿以补扇区和嵌入修复讨论奇点与宇宙起源：其数学内容是把当前表示中的不可延拓对象嵌入更丰富承载，并逐阶修复耦合障碍。物理非孤立由环境耦合判定，同调非孤立由补复形关系判定，两种判断分别给证书。宇宙论解释因此是附着在该形式模型上的条件化解释层，保持了生成链、时间性与宏微总同调的完整位置。

\section{开放系统的观测、领域迁移与受约束扩张}
\label{ss:open}

\subsection{统一的上下文修复引擎}

文献\cite{S25}把生命科学域 $V_2$ 和语义域 $V_3$ 统一为上下文索引的修复引擎族
\[
 S_v=(B_v,E_v\to B_v,\operatorname{Ext}_v,O_v,R_v,
      \Gamma_v^{\mathrm{adm}},\mathcal A_v,\mu_v,\Sigma_v,N_v),
 \qquad v\in\{2,3\}.
\]
其中上下文规定任务、证据与资源；纤维保存候选结构及修复层；$O_v$ 同时登记 Jacobi、Bianchi、引用、证书与尺度接续残差；$R_v$ 产生修复；允许路径类、作用量和统计测度执行选类；$\Sigma_v$ 验证并签发；$N_v$ 产生可复用正规形。经典强化学习或 SAC 可作为某个局部求解组件，而统一理论的组织中心是障碍、合法路径和证书。

对生命域，微观、介观、宏观、生命算子丛及生命态基底依次连接，附加算子包、约束、程序和证书。语义域使用同构的流程组织规范语义、上下文、证据及文本表达。这里可迁移的是任务接口及其交换关系，各领域的模型、数据与实验标准仍由专用插件提供。

开放性通过环境的非零耦合刻画：$D\Psi/dt=\mathcal U(\Psi)+\mathcal J(\Psi,e,t)$，其中 $\mathcal J\not\equiv0$。广义博弈记录子系统间的影响核，异构演化记录该网络上的实际轨道；它们是同一耦合对象的结构与过程两种读法。混沌、不确定性和有效开放性分别通过状态敏感性、观测纤维及隐藏扇区进入模型。能力扩张可以同时加强修复和影响其他扇区的能力，因此允许路径、用途约束、资源预算和证书必须随同扩张一并传递。

源稿以 GROMACS 分支说明外部计算核如何接入虚拟实验证书。其证书链依次保存源码与修改、构建和依赖、输入包、预处理、运行输出、分析、扩展接口和验证。完整性指规定实验类中的必要证据均可生成、导出和复核；计算核版本、分支新增能力与科学模型适用域分别标记。如此，领域仿真可以承担证据生产，主框架负责规定这些证据如何进入跨尺度判断。

\subsection{观测历史与终端纤维的充分性}

文献\cite{S28}将受限观测下的终端推断归为集合值逆问题。令 $\mathcal A$ 为可容许演化数据，$H_t:\mathcal A\to B_t$ 为观测历史映射，$\psi:\mathcal A\to Y$ 为终端量。能力元组
\[
 \Lambda_t=(\mathcal O_t,\mathcal M_t,\mathcal C_t,\mathcal L_t,\Pi_t)
\]
分别描述观测、模型、算力、时延和先验。相容条件筛出 $\mathcal A_t^{\Lambda}$，并定义
\[
 \Gamma_t^{\Lambda}(b)
 =\{\psi(a):a\in\mathcal A_t^{\Lambda},\ H_t(a)=b\}.
\]
该纤维保留相同历史下仍可能出现的终端结果。源稿中的“钱学森障碍”在本综述只沿用这一抽象含义：受限历史不能被稳定压缩为唯一终端结论。

\begin{theorem}[观测决定终端的充要条件]
\label{ss:observability}
在给定相容域内，存在满足 $\psi=f\circ H_t$ 的单值终端推断 $f$，当且仅当
\[
 H_t(a)=H_t(a')\ \Longrightarrow\ \psi(a)=\psi(a').
\]
若另有模连续函数 $\omega$ 使
\[
 d_Y(\psi(a),\psi(a'))\le\omega(d_B(H_t(a),H_t(a'))),
 \qquad \omega(r)\to0\ (r\downarrow0),
\]
则 $f$ 在可实现历史域上连续且具有同一稳定界。
\end{theorem}
\begin{certificate}
保留能力元组、可容许历史、纤维判据和稳定界。失败见证可由两个同历史而异终端的对象给出。选出纤维的一个连续截面，只证明存在连续选择；要声称终端已被唯一确定，还须验证纤维单点性。
\end{certificate}
\begin{proof}
必要性由 $\psi=fH_t$ 立即成立。充分性取 $f(b)=\psi(a)$，其中 $H_t(a)=b$；前件保证选不同代表不改变值。将代表代入模连续不等式即得 $d_Y(f(b),f(b'))\le\omega(d_B(b,b'))$。
\end{proof}

边缘算子可以改变观测分辨率、模型族、隐藏模式、更新时延或计算预算，因此改变终端纤维与可修复域。若真实对象始终被保留，且新证据使纤维嵌套收缩，最坏情形终端不确定性不会增大；若只是错误删去真实候选，则不构成有效修复。新增模式须与切换规则一并进入混合系统或微分包含。由此，结构完整性先决定求解的是何种问题，成本比较再决定怎样处理它。

稳健决策可写为 $\inf_u\sup_{a\in\mathcal A_t^{\Lambda}}J(u,a)$。更大预算只有在合法动作集或认证信息集按明确方式扩展时才给出单调性；单位预算收益、达阈时间和边际饱和均在同一目标及约束下比较。源稿所说完整性修复补齐当前任务的断裂链，完备性修复扩张可处理任务族；二者分别对应现有义务闭合与新义务覆盖。

\subsection{性变态射、力迫式扩张与有限交付}

文献\cite{S29}把性质变化的演化候选称为性变态射。候选的出现提供可研究的对象变化；要获得任务同一性，还需逐项给出编码、障碍、控制、有限截断和验证五类桥接义务。边缘算子负责修补接续，性质算子负责改变可用结构，两类算子组合产生有证书的扩张。

可用偏序 $p\preceq q$ 表示条件扩张，令每项待满足义务 $D_i$ 为稠密集合。若从任一 $p_n$ 都能有效选择 $p_{n+1}\succeq p_n$ 且 $p_{n+1}\in D_n$，则归纳构造的链逐项满足可数义务。极限对象是否存在以及是否保留旧解释，还须满足极限拼接和任务回读条件。该框架内的力迫属性由此具有构造内容；经典集合论 forcing 是在特定承载和投影下的严格实例。

无故障修复的世界观层主张、有限阶可证性和工程可交付性由三个不同证书连接。有限阶定理要求该阶修复算子实际存在并通过相容验证；有限交付要求当前任务在资源界内到达接受或显式边界。若一次交付必须实际展开无穷多层，则其工程流程未终止。相反，对持续增加的有限任务族，能够产生有限修复并保留既有证书，体现可扩展修复能力。这一定位保留高阶理论的作用，也使工程残差成为可定位、可处理的对象。

\section{Jacobiator-gap 动力、统计力学与方法空间}
\label{ss:dynamics}

\subsection{统一残差与可执行下降}

文献\cite{S17}把 Jacobiator-gap 从静态检查量转为修复动力。为统一后续求解，令 $z$ 包含数值参数、组织结构和证书状态，令
\[
 G(z)=(g_1(z),\ldots,g_m(z)),\qquad
 \Phi(z)=\frac12G(z)^{\mathsf T}WG(z),\qquad W>0.
\]
各 $g_i$ 分别对应任务残差、结构相容、证据缺口或约束违背。$\Phi$ 是选定度量下的势函数；它与具体效用的关系须通过任务映射建立。所谓“微观隧穿”在算法层指跨越当前表示的阻隔，命中另一个合法低残差代表，其可执行性由候选生成、合法性和下降证书共同给出。

\begin{lemma}[可执行方向的严格下降]
\label{ss:descent}
设 $\Phi$ 的梯度为 $L$-Lipschitz，实际可执行方向 $v$ 满足
\[
 \langle\nabla\Phi(z),v\rangle\le-a\|\nabla\Phi(z)\|^2,
 \qquad \|v\|\le b\|\nabla\Phi(z)\|,
 \qquad a,b>0.
\]
则当 $0<\eta<2a/(Lb^2)$ 时，
\[
 \Phi(z+\eta v)\le\Phi(z)
 -\left(a\eta-\frac{Lb^2\eta^2}{2}\right)\|\nabla\Phi(z)\|^2.
\]
\end{lemma}
\begin{certificate}
保存方向的实际实现、梯度误差、光滑常数、步长和约束检查。分段光滑或受约束问题改用相应有效方向与投影下降证书。
\end{certificate}
\begin{proof}
由光滑性，$\Phi(z+\eta v)\le\Phi(z)+\eta\langle\nabla\Phi(z),v\rangle+L\eta^2\|v\|^2/2$。代入两个界并合并即得；步长条件使括号为正。
\end{proof}

若进一步有 $\|\nabla\Phi\|^2\ge2\mu(\Phi-\Phi_*)$，且每步下降系数有正下界，则势能差以几何速率趋零。若零势能恰对应全部登记残差为零，则得到相应闭合。算子方向覆盖与该误差界承担从局部下降到全局结论的桥梁；组织同伦通过改变分组、层级和执行次序扩大可达方向，参数同伦在固定组织内优化已有方向。

\subsection{门控、积分触发、冷却与预算}

运行时将残差驱动速率写为 $r_t=r(G_t)$，以强度 $b_t$ 调用修复算子，并满足
\[
 r_t\,c(b_t)\le C_t.
\]
其中 $c(b_t)$ 为单次成本，$C_t$ 为可用算力。积分触发器积累 $r_t\Delta t$，达到阈值后执行，余量回存；安全门、冷却、阻尼、抗饱和与退化规则共同控制可调用区间。确定的输入、参数、并列规则与随机种子给出可重放的执行链。连续频率增加是否改善下降，要由单次修复效果、预算和饱和模型共同判断。

对开放扰动，可将一次总变化分解为
\[
 \Phi(z_{t+1})-\Phi(z_t)
 =\Delta_t^{\mathrm{repair}}+\Delta_t^{\mathrm{environment}}.
\]
若修复项被证明为负，仍需记录环境注入才能解释总残差轨道。目标效用采用独立量 $U$，净收益由其实际兑现与成本决定。于是门控稳定、残差下降和最终收益分别有各自证书，避免把机制改进重复计入应用收益。

\subsection{PDE 策略簇与软件层统计结构}

文献\cite{S32}将求解方法本身参数化。策略态
\[
 \theta=(\text{离散尺度},\text{网格},\text{边界处理},\text{时间推进},
          \text{迭代器},\text{正则项},\text{停止规则})
\]
组成格点或带邻接结构的方法空间。对策略的残差、误差、费用和证书缺口定义能量，方法选择便成为该空间中的分布演化与合法迁移。连续可微的参数使用方向导数，离散方法切换使用显式邻接差分，二者通过统一评价接口接入同伦扭曲。

将命题~\ref{fa:prop-gibbs}应用于策略态和策略能量，即得有限策略族的 Gibbs 权重。
若所有非最优态与最优值间隔至少 $\delta>0$，则非最优态的总概率被常数倍 $e^{-\beta\delta}$ 控制。这一结论描述分布质量的集中；采样能否到达该分布，由支持、混合和调度证书确定。频域或结构域分析进一步识别慢模态、相关方向和可命中通道，使策略选择使用整个轨迹的结构信息。

路径权重、干涉式表示、低作用量通道与末端测量构成原文所说近似量子软件层的结构接口。它们在经典软件上组织统计搜索，物理量子实现和复杂度结论需要各自数据。自动证明、科研与设计则在具备可计算残差、候选策略及独立验证器时接入同一策略簇模板；“高阶泛函约束”因此是任务统一的承载语言，具体可解性仍由任务接口提供。

\section{LHTS、TPH 与高阶同伦命中}
\label{ss:hitting}

\subsection{分量残差、延迟信用与 D 结构反拉回}

LHTS 即格点同伦扭曲评分体系，负责在给定问题空间内产生可执行运动；TPH 即轨迹函数参数同伦控制，负责比较和更新生成这些运动的轨迹函数。两者分别对应运动内核与外层选择控制，组成文献\cite{S34,S35}的统一接口。

设 $T_i^\alpha$ 为单格点扭曲，记录分量效果
\[
 \delta_{ij}(z;\alpha)=g_j(z)-g_j(T_i^\alpha z),
 \qquad \boldsymbol\delta_i=(\delta_{i1},\ldots,\delta_{im}).
\]
当效果延迟显现时，信用依据完整回放归入原动作，并写入持久评分张量。分量记录保留不同约束的损益，标量评分只作任务权重下的投影。尤其当某些约束不可补偿时，应先检查逐分量门槛，再在通过者中比较总分。

\begin{proposition}[压力倒决策的雅可比反拉回]
\label{ss:pressure}
令当前残差压力为 $q\in\mathbb R^m$，微小动作参数为 $a\in\mathbb R^r$，且 $G(z(a))$ 在 $a=0$ 可微，雅可比为 $J$。则一阶加权残差变化为
\[
 q^{\mathsf T}\bigl(G(z(a))-G(z(0))\bigr)
 =(J^{\mathsf T}q)^{\mathsf T}a+o(\|a\|).
\]
因此 $J^{\mathsf T}q$ 给出动作空间的一阶压力评分；具体选择再结合可行方向、控制幅度与费用。
\end{proposition}
\begin{certificate}
保存残差分量、压力生成规则、实际动作顺序、雅可比或有限差分估计，以及延迟信用的回放标识。
\end{certificate}
\begin{proof}
将可微展开 $G(z(a))=G(z(0))+Ja+o(\|a\|)$ 左乘 $q^{\mathsf T}$，并用内积的转置恒等式即得。
\end{proof}

组织动作也能进入该格点体系：实单、虚单、虚仓的组织变换，或抽象任务中的槽位、分组、顺序和层级调整，都可视为有类型的生成元。由于动作通常不交换，候选应保存有序词 $a_1\cdots a_k$、合法重写规则及正规形。Parikh 计数只记录各生成元出现次数；若 $ab\ne ba$，同计数的动作词可有不同终点与费用，因此计数频域的等价不能替代执行路径等价。

LHTS 的完整性接口包括三项：有限工程候选的覆盖或可靠排除，高阶几何与极限保真，以及不被评分机制消灭的探索支持。它们分别服务于工程全局比较、连续极限提升和持续探索。分支定界、动态规划或穷举可提供第一项，不要求所有任务采用同一种求解器。相应子簇包括 D-LHTS（分量雅可比压力）、AJ-LHTS（退火跳跃）、SJ-LHTS（评分跳跃）、MT-LHTS（隧穿与避障）、EG-LHTS（工程全局覆盖）及 $L_\infty$-LHTS（高阶闭合与极限提升）；其差别由候选生成和证书接口确定。

组织同伦的应用由固定工作域给出。以 OpenRA 沙盘为例，将运行态写成组织 $\Sigma$ 与部署 $D$；若所有合法动作的前提和结果均由 $(\Sigma,D)$ 决定，且每个动作都可分解为合法重组织与重部署，则两种表达具有相同有限执行。证明按动作词长度归纳，并逐步替换生成元；若只保留端点而丢失中间费用，便须另补路径证书。在固定槽位扩张中，库存变化可由槽位类型 $\tau(u)$ 的空、非空和换型表示，每次转换仍逐项核验资源与费用约束。量化域的实单--虚单--虚仓使用相同组织投影，并保留账本、敞口及费用。两类应用展示共同组织内核在不同任务解释下的实例化\cite{S17}。

\subsection{首次命中、概率路径与幂化关系}

文献\cite{S33}将目标设为低残差集 $B_\epsilon=\{z:E(z)\le\epsilon\}$。轨迹函数 $\Gamma$ 诱导路径测度 $\mathbb P_\Gamma$，定义
\[
 \tau_\epsilon=\inf\{t:z_t\in B_\epsilon\},\quad
 J_T(\Gamma)=\mathbb P_\Gamma(\tau_\epsilon\le T),\quad
 Q_T=1-J_T.
\]
若目标由频谱窗口、相位或模态关系定义，事件依赖整个有限路径 $z_{0:T}$，仍由相应可测路径空间承载。确定性轨迹是 Dirac 测度的特例，随机样本与确定性的选择规则可以同时存在。

\begin{theorem}[条件命中与生存概率]
\label{ss:survival}
在生存概率为正时，设 $p_t$ 为前 $t-1$ 次均未命中条件下，第 $t$ 次命中的条件概率；一旦生存概率为零，后续约定 $p_t=0$ 并保持 $Q=0$。则
\[
 Q_N=\prod_{t=1}^N(1-p_t)
 \le\exp\left(-\sum_{t=1}^Np_t\right).
\]
若 $\sum_t p_t=\infty$，则最终命中的概率为一。若每次 $p_t\ge p>0$，失败概率至多为 $e^{-Np}$。
\end{theorem}
\begin{certificate}
记录目标集、路径法律、条件概率下界及接受判据。条件概率取在未命中历史事件上；若使用逐历史下界，则将其通过条件期望传给 $p_t$。这允许自适应采样而无需假设独立。
\end{certificate}
\begin{proof}
由条件概率乘法公式递推得到有限乘积。用 $1-x\le e^{-x}$ 得指数界；概率测度对递减未命中事件连续，故级数发散时 $Q_\infty=0$。
\end{proof}

有限时域算法与理想分布的联系由混合证书量化。若运行 $b$ 步后的样本分布 $\mu_b$ 满足 $\|\mu_b-\pi\|_{\mathrm{TV}}\le\delta_b$，对 $m$ 次独立重启的终端样本，至少一次落入目标集 $B$ 的概率满足
\[
 \bigl|J_{m,b}-[1-(1-\pi(B))^m]\bigr|\le m\delta_b.
\]
单次事件概率之差至多为 $\delta_b$，而 $x\mapsto1-(1-x)^m$ 在 $[0,1]$ 上为 $m$-Lipschitz，故得此界。若目标计入中间时刻的首次命中，则把整段路径作为样本，对路径分布验证同一全变差界；样本估计误差另行加总。这样，实际时域运行与分布命中质量的优选通过明确误差联系。

当 $0<Q_T<1$，定义命中指数 $C_T=-\log Q_T$。任意增益 $\Delta C_T\ge0$ 给出
\[
 Q_T^{\mathrm{new}}=Q_T^{\mathrm{old}}e^{-\Delta C_T}.
\]
若 $C_T^{\mathrm{new}}=aC_T^{\mathrm{old}}$，则 $Q_T^{\mathrm{new}}=(Q_T^{\mathrm{old}})^a$。原文的“幂化命中”因此是命中指数改善的精确重写；幂指数的增益仍需由具体扭曲计算。薄通道、障碍两侧的合法代表和低残差区域可以通过同伦方式组织，但通道的几何可达性与测度上的命中质量各有其证书。

正厚度通道与芯线的关系可以精确保存。设 $T=D^k\times[0,1]$，$C$ 为两端盘与中间芯线之并，若给定相对两端盘的强变形收缩 $T\searrow C$，则在同一端口映射下将它们粘到原空间 $X_0$，收缩仍能下降到粘合空间，得到 $X_0\cup T\simeq X_0\cup C$。证书包含端口映射、收缩和端口固定条件；证明是在 $X_0$ 上取恒等、在 $T$ 上取所给收缩，再由粘合相容性得到整体同伦。由此，通道厚度可在保持该同伦任务时简化，路径的概率与费用仍由各自结构计算。

\subsection{正交命中与高阶雅可比间隙}

在有限空间中取指数扭曲 $\mu_{\eta,v}(x)=\mu(x)e^{\eta v(x)}/Z(\eta)$，直接求导得
\[
 \left.\frac{d}{d\eta}\mathbb E_{\mu_{\eta,v}}f\right|_{\eta=0}
 =\operatorname{Cov}_\mu(f,v).
\]
因此与受保护目标 $u$ 正交，即 $\operatorname{Cov}_\mu(u,v)=0$，保持其期望的一阶不变；命中是否加速由 $\operatorname{Cov}_\mu(1_B,v)>0$ 决定。高阶正交进一步在已使用方向的正交补中寻找剩余命中梯度，使各阶新增方向对应尚未修复的缺口。

\begin{theorem}[高阶正交方向的有效性证书]
\label{ss:orthogonal}
令 $g_k=\nabla C_T(\Gamma^{[k-1]})$，$W_{k-1}$ 为已用方向空间，并令
\[
 b_k=\operatorname{proj}_{W_{k-1}^{\perp}}g_k\ne0.
\]
若实际可执行方向 $\delta_k=J_kd_k$ 满足
\[
 \|\delta_k-b_k\|<\frac{\|b_k\|^2}{\|g_k\|},
\]
则 $DC_T[\delta_k]>0$，足够小的正步长提高命中指数并降低未命中率。
\end{theorem}
\begin{certificate}
记录 $W_{k-1}$、梯度、实现雅可比 $J_k$、控制 $d_k$ 和误差界。可用
\[
 \mathfrak G_k=\frac12\|J_kd_k-b_k\|^2
 +\frac{\rho_k}{2}\|\operatorname{proj}_{W_{k-1}}J_kd_k\|^2
\]
同时度量方向逼近与正交泄漏；具体步长另过合法性与增益验证。
\end{certificate}
\begin{proof}
设 $e_k=\delta_k-b_k$。正交投影给出 $\langle g_k,b_k\rangle=\|b_k\|^2$，故
\[
 DC_T[\delta_k]
 =\|b_k\|^2+\langle g_k,e_k\rangle
 \ge\|b_k\|^2-\|g_k\|\|e_k\|>0.
\]
可微性保证足够小正步长有真实增益，再用 $Q_T=e^{-C_T}$ 得未命中率严格下降。
\end{proof}

这一公式解释 Jacobiator-gap 如何指引更高阶命中：先压低实际动作与理想修复方向之间的间隙，再由经过验证的方向转为命中指数增长。各阶优化无需每次重新解决整个原始问题；它处理的是尚未被既有方向覆盖的剩余分量。

\subsection{TPH 选择、递归回填与确定性控制}

TPH 从 LHTS 数据 $D_t$ 生成有限非空候选族 $\{\Gamma_t^{(r)}:r\in R_t\}$，以命中概率、预期残差下降、费用、方差及方向覆盖误差构成评分 $\mathcal L_D$。固定并列规则得到 $r_t^*$，执行所选轨迹并将实际结果回填至评分与候选生成器。高阶候选可以同时携带概率路径同伦类、递归轨迹族同伦类和具体可执行代表：
\[
 \gamma_D=([\mathbb P_\Gamma]_{\mathrm{prob}},
           [\mathcal T_\infty]_{\mathrm{rec}},\Gamma).
\]
其中 $\mathcal T_{k+1}=\mathfrak F(\mathcal T_k)$ 表示生成、筛选、回填的递归，满足相应尺度兼容时进入广义分形子类。评分与最大元选择对应 D-Decision，递归深度与新一轮修复对应 D-Depth。

\begin{theorem}[预测裕量与有限可执行闭环]
\label{ss:tph}
设每轮候选有限非空，评分及验收判定全定义，并列规则确定，接受的更新保持合法状态域。候选未通过验收时按确定次序继续检查，全部未通过时返回未闭合义务并停止。设待接受候选预测下一步势能为 $\widehat\Phi_{t+1}$，误差界为
\[
 |\Phi(D_{t+1})-\widehat\Phi_{t+1}|\le e_t.
\]
若接受条件为
\[
 \Phi(D_t)-\widehat\Phi_{t+1}\ge e_t+m_t,\qquad m_t\ge0,
\]
则每轮都有确定的接受或停止结果，且任何已接受更新满足 $\Phi(D_{t+1})\le\Phi(D_t)-m_t$。若 $\Phi\ge0$，无限已接受更新链上的势能值收敛；若该链无限延续，且对每个 $\epsilon>0$，残差高于 $\epsilon$ 时始终选中至少下降 $c_\epsilon>0$ 的合法候选，则残差趋零。
\end{theorem}
\begin{certificate}
保存候选集合、评分、并列次序、预测误差、真实回放、合法更新和下降裕量。随机轨迹的种子或样本记录也是回放输入；确定性指给定这些输入后选择与验收可重现。
\end{certificate}
\begin{proof}
有限非空集合上的全定义评分及并列规则给出确定检查次序；有限全定义验收产生一个被接受候选或明确停止结果。合法更新使已接受前缀可逐步归纳。误差界与接受条件合并得到下降式。无限已接受链上的非负单调势能序列必有极限。若其残差无穷多次高于某个 $\epsilon$，每次至少下降 $c_\epsilon$，将使非负势能降为负，矛盾，故残差趋零。
\end{proof}

高阶方向覆盖的作用，是提供定理中持续可用的下降候选。单纯的“每次略有下降”不足以给出残差归零，因此记录离零区域的定量下降或等价梯度误差界。若只证明经验目标驻点收敛，就签发该层结论；若再有样本一致收敛、离散化一致性、可行域紧性和无伪驻点条件，则经验选择可以提升为理论选择，并获得极限解。这把有限样本的可回放执行、概率一命中和数学极限归零接在同一条明确的证书链上。

\section{PDE 的测度表达、命中求解与解释等价}
\label{ss:pde}

\subsection{由局部微分到 L-积分残差}

文献\cite{S36}的集中贡献，是把偏微分方程的表达、求解和解释分为三个相接的接口。设任务为
\[
 \mathcal P[u]=0,\qquad \mathcal B[u]=0,
\]
其中 $\mathcal B$ 合并边界及初值条件。法空间测度结构记为
\[
 \mathfrak L=(X,\Sigma_L,\mu_L,\nabla_L,\mathcal H_L,\Pi_L,\mathcal C_L).
\]
当变化测度 $D_Lu\ll\mu_L$ 时，在适用的 Radon--Nikodym 条件下定义 $\partial_Lu=dD_Lu/d\mu_L$。此时
\[
 D_Lu(A)=\int_A\partial_Lu\,d\mu_L,
\]
微分成为变化测度的局部密度，积分保存整体变化。本文使用的 $L$ 是法空间测度及语义分辨率记号，与前述 $\mathsf L_i$ 修复层号分别登记。

令测试对空间为 $\mathcal T=\mathcal T_X\times\mathcal T_{\partial X}$，定义
\[
 W_L(u;\varphi,\psi)
 =\int_X\langle\mathcal P[u],\varphi\rangle\,d\mu_L
 +\int_{\partial X}\langle\mathcal B[u],\psi\rangle\,d\sigma_L,
\]
以及
\[
 E_L(u)=\sup_{\|(\varphi,\psi)\|_{\mathcal T}\le1}
              |W_L(u;\varphi,\psi)|^2.
\]
若原问题采用分布导数，则先由积分分部给出适用的弱表达，并把函数空间、迹和测试条件列入任务定义。

\begin{theorem}[表达、求解与解释的联合迁移]
\label{ss:pde-equivalence}
假定测试族分离规定残差，路径表示空间 $\Gamma_L$ 有投影 $\Pi_L:\Gamma_L\to\mathcal U$，并满足：
\begin{enumerate}
 \item 每个原问题解 $u$ 均有合法路径代表 $\gamma$，使 $\Pi_L\gamma=u$；
 \item 接受路径须验证 $E_L(\Pi_L\gamma)=0$，或明确的工程阈值；
 \item 规定解释交换：$\mathcal C_L^{\mathrm{path}}(\gamma)=\mathcal C_L^{\mathrm{PDE}}(\Pi_L\gamma)$。
\end{enumerate}
则零残差层上
\[
 \mathcal P[u]=0\wedge\mathcal B[u]=0
 \quad\Longleftrightarrow\quad E_L(u)=0,
\]
原解与合法零残差路径在投影下互相覆盖，并对规定任务解释等价。由此，计算中心可以迁移至路径空间中的低残差命中。
\end{theorem}
\begin{certificate}
保存函数与测试空间、残差分离性、边界迹、路径覆盖、投影、接受证书和解释交换图。零残差按所选解概念理解：通常为几乎处处或弱意义成立；升级为经典逐点解还需相应正则性。
\end{certificate}
\begin{proof}
原残差为零时所有测试值为零，故 $E_L=0$。反向若 $E_L=0$，利用线性缩放可得全部测试值为零；分别令一个测试分量为零，再用分离性推出内部和边界残差均为零。覆盖条件把每个原解送入合法路径表示，零残差接受分支把每个相应路径送回原解；解释交换保证两种表示给出同一任务意义。三项合取完成表达、求解对象与解释的迁移。
\end{proof}

\subsection{路径测度与独立接受证书}

在合法路径族上定义作用量 $\mathcal A_L$ 与参考测度 $M_L$，若 $0<Z_\beta<\infty$，则
\[
 \nu_\beta(d\gamma)=Z_\beta^{-1}e^{-\beta\mathcal A_L(\gamma)}M_L(d\gamma).
\]
作用量可组合残差、Jacobiator-gap、费用、风险与任务效用，用于引导候选。最终接受仍独立检查投影残差和硬约束，因为加权目标之间可能存在补偿。LHTS 提供路径的局部动作与信用，TPH 调整轨迹函数，高阶正交搜索未覆盖方向；三者共同把原来的显式偏微控制组织为隐式测度表达下的可验证路径求解。

若旧算法的合法轨迹可嵌入新路径族，且费用与解释一同保持，则旧算法成为新方法族中的候选实例；扩充路径族能够增加可比较的合法方案。这一包含关系给出方法空间的结构增益。其实际速度、样本效率和净收益，由候选生成、验证、回落及总成本共同决定，按相同任务记录。

\subsection{工程阈值、严格极限与兑现归因}

有限输出 $u_n$ 若满足 $E_L(u_n)\le\epsilon_n$，得到可复核的低残差代表。若同时有紧性产生收敛子列 $u_{n_j}\to u$，$\epsilon_{n_j}\to0$，且 $E_L$ 下半连续，则
\[
 0\le E_L(u)\le\liminf_{j\to\infty}E_L(u_{n_j})=0,
\]
从而由定理~\ref{ss:pde-equivalence} 获得原问题解。若还需误差估计，则另给稳定界
\[
 \operatorname{dist}(u_n,\operatorname{Sol})
 \le C E_L(u_n)^\alpha.
\]
残差、误差和解存在性由这些明确接口相接。问题转化本身一旦建立，就可使原 PDE 保留为解释与校验依据，而求解器采用等价的测度和命中表示。

文献\cite{S36}还将收益落实到有限回放域。对每步有符号贡献 $I_t=I_t^+-I_t^-$，恒有
\[
 \sum_{t<T}I_t=\sum_{t<T}I_t^+-\sum_{t<T}I_t^-.
\]
净贡献为正可等价写成累计正贡献大于累计负贡献，并不要求每个时刻为正。要把残差下降归因于某个修复机制，需要同域输入、对照、费用及排他项的回放证书；文件中给出的证书模板与实际已经填入数据的事件记录分别保存。

最终，这一求解谱系将\emph{对象表示、障碍识别、合法扩张、路径选择、残差验收与解释回读}连成一个流程。表达可以变化，任务意义由交换图保持；局部失败可以进入更高层修复，新增自由度通过证书回接；有限工程结果与严格极限则各自有可复验的完成条件。这为后续“率与商”的哲学收束提供了具体算子、路径及成本对象。


\section{基础障碍谱系：从问题名录到同一证书结构}
\label{pc:spectrum}

基础障碍谱系把集合与逻辑、几何与数论、物理与计算、生命与认知中的困难组织成同一种可比较对象。其共同项包括：问题需要保留哪些谓词，现有表示在何处失配，失配能否构成明确复形中的障碍，允许增加何种见证，终端输出如何回到原任务。由此，问题之间的统一落实到\emph{定义域、结构映射、障碍对象、修复操作与接受条件}，而各问题自身的对象和结论继续保留。文献~\cite{S37} 给出跨学科谱系，文献~\cite{S38} 进一步将这一谱系的数学接口展开为任务商、微分交换和合法回读。

\subsection{证明增强、构造定理与原子化审计}

\begin{definition}[问题对象与七槽证书]
把问题 $Q$ 登记为
\[
 \mathcal O_Q=(Q,\mathcal D_Q,\mathcal P_Q,\mathcal B_Q,
              \mathcal R_Q,\mathcal A_Q),\qquad
 c_Q=(N,G,P,C,A,B,T),
\]
其中 $\mathcal D_Q$ 是对象域，$\mathcal P_Q$ 是必须保持的谓词，$\mathcal B_Q$ 是已定义的边界或障碍，$\mathcal R_Q$ 是允许修复族，$\mathcal A_Q$ 是接受条件。证书七槽依次记录对象名、生成数据、谓词、上下文约束、审计、边界障碍及终端判断。接受是各项实际证明义务的合取；必需见证未完成时，保留相应前提或未决状态。
\end{definition}

在这一表述中，命题之间的逻辑蕴含与操作引起的因果变化分别记录。前者由合法推理规则连接，后者需要结构方程、作用对象和外生条件固定的比较分支。对同一记号 $d$ 也按类型区分：只有给出分次空间及 $d^2=0$ 的实现才称为余链微分；规范化、验证或状态更新采用自身的算子记号。由此避免把操作名称相似当作结构等同。

验证结构可由集合、函数、有限谓词、残差队列和接受集合构造；应用名称只是对象标签，真正进入推理的是具体操作数、条件与见证。语义信息可统一记录为有效语义、保留不变量、来源证书和待修复残差的组合，概率熵则描述指定概率分布的另一类量。吸收、剪枝和回写为处理过程规定有向转移；这种操作上的时间方向及其与物理过程的关系，各由转移和桥接证书确定。

\begin{theorem}[证书增强的忠实性与相对保守性]
\label{pc:conservative}
设原理论 $\mathcal T$ 的公式被增强为
\[
 \iota(F)=(F,\operatorname{NF}(F),\operatorname{Inv}(F),
                  \operatorname{Cert}(F),\operatorname{Audit}(F)),
 \qquad \pi\iota(F)=F.
\]
若新增规则仅作定义展开、证书登记与已证规则的审计，并且遗忘投影把每条新增规则变为原理论中的合法步骤、恒等式或空步骤，则原证明可以逐步增强；任一增强证明投影后仍为原理论的证明。
\end{theorem}
\begin{certificate}[证明长度归纳]
保存有限证明序列、每步父结论及规则实例；增强时保留原公式核，遗忘时逐条登记规则的投影。归纳谓词为“长度不超过 $n$ 的前缀及其投影均保真”。
\end{certificate}
\begin{proof}
长度零没有待验证结论；初始公式保留原理论的基础见证。设前 $n$ 步已保持。原规则只使用此前结论，保持公式核即可实施同一推理；新增审计规则只增加已定义数据，其投影按前提为合法步骤或可删除步骤。因此第 $n+1$ 步仍保真。有限归纳得到双向结论。该证明同时说明：系统特有的闭合规则可以在基础对象和构造子的保持证明之上成为定理；基础语言本身及尚待实例化的可容许性、收敛性条件仍须明确声明。
\end{proof}

“原子推理步”相对于声明的语义签名确定。复合证明先展开为有类型的有限依赖图，再收束到变量读取、常数、单个矩阵条目、基础运算和已核对的规则实例。每个结点保留未解除假设集合
\[
 \mathsf{Open}(v)=\bigcup_{u\in\operatorname{Par}(v)}\mathsf{Open}(u)
                   \cup\mathsf{NewPrem}(v),
\]
并仅按明确规则解除。拓扑顺序归纳保证已接受结论可追溯到来源叶；事件接受要求必需计算已完成且开放假设已经按任务要求处理。无限命题可以由有限的带参数符号证明验证，其证明对象与无限枚举分开。

\subsection{障碍谱系的完整压缩}

表~\ref{pc:hilbert-table} 将希尔伯特二十三问题按原文五类聚类及其后补充的第5题合并，形成不重复且完整的索引分区。这里“证书化化解”统一指向规定任务的规范化、见证与接受链；若要推出原问题的特定结论，还须在同一对象域内提交该结论的回读及等价证明。

\begin{table}[htbp]
\centering\small
\caption{希尔伯特问题谱系的合并索引与证书接口}
\label{pc:hilbert-table}
\begin{tabular}{p{0.15\textwidth}p{0.23\textwidth}p{0.50\textwidth}}
\hline
聚类 & 问题编号 & 统一登记的结构义务\\\hline
逻辑与相容性 & 1、2、10 & 连续统分层、形式相容性、离散搜索及失败分支。\\
连续群结构 & 5 & 连续性与局部可微表示之间的承载和映射条件。\\
物理公理化 & 6 & 跨物理扇区的共同对象、量纲、约束与接受映射。\\
PDE与变分 & 19、20、21、22、23 & 正则性、边值、单值化、变分、奇点及可回读残差。\\
几何与拓扑 & 3、4、15、16、18 & 几何基础、交点、空间类型、铺砌和同调不变量。\\
数论与解析 & 7、8、9、11、12、13、14、17 & 函数零点、算术结构、代数不变量、有限生成及锚定证书。\\\hline
\end{tabular}
\end{table}

六行编号两两不交，数量为 $3+1+1+5+5+8=23$，并集为 $\{1,\ldots,23\}$。这一有限覆盖只负责组织问题；每项证明仍由该项的输入与映射承担。扩展谱系进一步归纳为表~\ref{pc:extended-table}，同一问题在不同章节中出现的对象化、桥接和总证书合并到同一行。

\begin{table}[htbp]
\centering\small
\caption{扩展基础障碍谱系及其共同接口}
\label{pc:extended-table}
\begin{tabular}{p{0.34\textwidth}p{0.54\textwidth}}
\hline
问题族 & 在统一框架中的组织方向\\\hline
Hodge、BSD、Yang--Mills质量缺口、Navier--Stokes & 同调类与代数代表、零点阶与残差维度、可证自由度及正下界、PDE残差与全局回读；各桥接分别提供其对象与见证。\\
Gödel、图灵停机、算法信息边界 & 区分原理论命题、给定证书的可验证性、搜索与执行；以有效证明、边界见证、证书深度及未决项组织接口。\\
量子引力、测量与黑洞信息 & 离散连续承载的相容、终端观察、投影与总同调信息保持。\\
Levinthal、复杂多体、微扰失效、各态历经破缺 & 从庞大候选空间提取任务相关轨道、不变量、边界和命中条件；给出求解及证据的实际访问成本。\\
Shannon语义盲区、Loschmidt时间箭头 & 在概率信息之外登记语义内容、保留类与修复残差；明确吸收、剪枝和回写形成的有向过程。\\
Wigner有效性、复杂涌现 & 数学表示与物理表示的共同正规型、任务商与端到端宏观证书。\\
Arrow社会选择 & 多偏好约束、冲突残差、允许的表示或规则扩张及其重新认证。\\
Langlands、符号接地与约束闭合 & 离散连续对象间的结构字典，符号对真实来源的锚定，有限谓词及高阶相容约束的验证。\\
意识硬问题与G-Qualia & 以残差识别、回写、收紧、语义重新确认和审计构造内部自修复轨道模型。\\\hline
\end{tabular}
\end{table}

\begin{proposition}[有限谱系证书的合成]
设有限问题集 $I$ 已被若干聚类覆盖，且每个 $i\in I$ 都有定义域相同、结论明确的证书 $c_i$。则可构造总证书 $c_I=(I,\kappa,(c_i)_{i\in I})$，验证全部规定结论。
\end{proposition}
\begin{certificate}[覆盖与逐项验证]
保存无遗漏索引表、聚类映射及各项有效证书；验证谓词为 $P(n)=\bigwedge_{i\le n}\operatorname{Verify}(c_i)$。
\end{certificate}
\begin{proof}
空合取成立；第 $n+1$ 项由其自身证书验证，与已通过前缀合取得到 $P(n+1)$。有限归纳完成。覆盖证明保证每项恰有归属，不代替任意单项所需的数学构造。
\end{proof}

\subsection{共同正规型与数学有效性的桥接}

\begin{theorem}[共同任务表示的解释与预测保持]
\label{pc:wigner}
设数学实现 $M$ 与观测实现 $X$ 有到共同任务域 $Q$ 的映射 $q_M,q_X$，规定观察满足 $F_M=\bar Fq_M$、$F_X=\bar Fq_X$。若初态同像，允许作用及全部响应在两侧对应，且后继与 $q_M,q_X$ 交换，则任何有限对应执行的规定观察相同。
\end{theorem}
\begin{certificate}[共同正规型的实质内容]
保存输入来源、状态对应、观察因子化、动作响应对应、后继交换以及回读；对物理任务另保存观测误差、量纲和模型适用域。
\end{certificate}
\begin{proof}
初态同像使两侧观察均为同一个 $\bar F(q_Mm)=\bar F(q_Xx)$。设第 $n$ 步同像，后继交换把对应动作响应送到共同后继，故第 $n+1$ 步同像。按长度归纳得到全部有限执行上的观察等值。近似实现将等值替换为后文的误差递推。
\end{proof}

该定理给出了数学为何能够有效描述某类过程的结构解释：共同保持的关系与不变量，使计算结果具有可检验的回读。数学表达、真实观测与它们之间的映射分别承担证明义务。对涌现而言，如果宏观任务已通过可访问不变量和边界证书因子化，中间帧可以从\emph{该任务的输出义务}中省去；产生这些不变量和证书的工作仍须计量。对意识而言，G-Qualia在此承担可审计自修复轨道的形式模型角色，经验体验与该模型之间的对应由另行声明的观测与解释接口承接。这保留了谱系的统一解释方向，也把每条桥接落实到能够检查的数学对象。

\section{任务充分商：同一性的构造、执行保持与必要细化}
\label{pc:quotient}

\subsection{任务同一性与性质塔}

\begin{definition}[任务结构及精确任务商]
令
\[
 \mathcal M=(X,A,\Xi,\delta,U,G,\mathcal T),
\]
其中 $X$ 是状态集，$A(x)$ 是合法动作，$\Xi(x,a)$ 是非空响应集，$\delta(x,a,\xi)$ 是后继，$U$ 是规定观察，$G$ 是目标集，$\mathcal T$ 是必须保持的测试族。采用共同动作响应标签时，满射 $q:X\to\bar X$ 是精确任务商，若
\begin{equation}
\label{pc:eq-exact}
\begin{aligned}
 A(x)&=\bar A(qx),& \Xi(x,a)&=\bar\Xi(qx,a),\\
 q\delta(x,a,\xi)&=\bar\delta(qx,a,\xi),& G&=q^{-1}(\bar G),\\
 U(x,a,\xi)&=\bar U(qx,a,\xi),& f&=\bar f\circ q_f\quad(f\in\mathcal T).
\end{aligned}
\end{equation}
这里 $q_f$ 是状态、动作响应记录或有限执行上下文上的诱导映射；动作改名时另给覆盖全部合法动作和响应的双向对应。
\end{definition}

任务同一性由必须保留的测试确定。对状态测试族 $\mathcal P$，令
\[
 x\sim_{\mathcal P}y\quad\Longleftrightarrow\quad
 \forall f\in\mathcal P, f(x)=f(y).
\]
新增测试产生更细的任务关系；层间唯一粗化映射及复合保持由命题~\ref{fa:prop-propertytower}给出，已有观察沿粗化继续可读。动态任务还必须把未来合法行为纳入测试，不能仅以当前效用值决定等价。

\begin{lemma}[观察因子化的充要条件]
\label{pc:factorization}
给定满射 $q:X\to Q$ 和观察 $F:X\to Y$，存在唯一 $\bar F:Q\to Y$ 使 $F=\bar Fq$，当且仅当 $F$ 在每个 $q$ 纤维上恒定。
\end{lemma}
\begin{certificate}[同像比较]
逐纤维验证 $q(x)=q(y)\Rightarrow F(x)=F(y)$，或提交一对同像异值输入作为否决见证。
\end{certificate}
\begin{proof}
若因子化成立，同像输入经 $\bar F$ 得到同值。反向定义 $\bar F(qx)=F(x)$，纤维恒定保证代表元无关，满射保证定义域完备和唯一性。
\end{proof}

\subsection{最粗动态充分商与有界策略}

\begin{theorem}[有限任务的演化饱和]
\label{pc:saturation}
设状态域 $X$ 与作用响应标签集 $L$ 有限，把未定义后继送入新状态 $\bot$，并使全部扩展后继 $\widehat\Phi_\ell$ 固定 $\bot$。令有限初始测试族 $\mathcal P_0$ 包含目标、观察、合法性、响应和 $\mathbf1_{\{\bot\}}$，递归置
\[
 \mathcal P_{n+1}=\mathcal P_n\cup
 \{f\circ\widehat\Phi_\ell:f\in\mathcal P_n,\ell\in L\}.
\]
若初始分区有 $b_0$ 块，诱导分区至多严格细化 $|X\sqcup\{\bot\}|-b_0$ 次即稳定。稳定商保持所有有限标签路径后的初始测试，是保持这些测试与后继的最粗商。
\end{theorem}
\begin{certificate}[分块与稳定性]
每次细化记录同块异值对及对应的观察路径；稳定时保存每个块在每个标签下的后继块及初始测试常值。
\end{certificate}
\begin{proof}
严格细化至少增加一块，块数有给定有限上界。某轮不再细化时，$x\sim y$ 蕴含 $f(\widehat\Phi_\ell x)=f(\widehat\Phi_\ell y)$ 对当前全部测试成立，故后继仍同类。按路径长度归纳，全部有限后继测试均保持，此后不再细化。反之，任何保持初始测试和后继的等价关系，都由同一归纳保持所有生成测试，因而只能比稳定关系更细。定义域标记和动作响应测试使稳定商满足式~\eqref{pc:eq-exact}。
\end{proof}

\begin{theorem}[鲁棒可达性与策略回落]
\label{pc:reach}
定义
\[
 W_0=G,\qquad
 W_{n+1}=W_n\cup\{x:\exists a\in A(x)\ \forall\xi\in\Xi(x,a),\ 
                         \delta(x,a,\xi)\in W_n\}.
\]
精确任务商满足 $W_n=q^{-1}(\bar W_n)$。商层的秩下降策略可合法回落，并在规定步数内对全部允许响应到达原目标。
\end{theorem}
\begin{certificate}[目标、量词域与秩]
保存目标原像式、共同动作响应域、后继交换及
$\rho(\bar x)=\min\{n:\bar x\in\bar W_n\}$；在正秩状态保存一个使所有后继秩下降的动作。
\end{certificate}
\begin{proof}
零步是目标原像式。设结论在 $n$ 成立，动作和响应的量词域由精确商相同，后继交换将 $\delta(x,a,\xi)\in q^{-1}(\bar W_n)$ 等价变为 $\bar\delta(qx,a,\xi)\in\bar W_n$。原像保持并集，得到 $n+1$ 步。回落动作在原层合法，每一步把非负整数秩至少减少一；至多初始秩步后进入目标。该证明同时保持任意有限执行的观察和接受值。
\end{proof}

有限完整转移表进一步给出可构造求解器：若核心有 $n$ 个状态、表大小为 $B$，逐轮扫描鲁棒前驱可在 $O(nB)$ 预处理中生成胜域与秩策略；固定点补集中，每个合法动作都有至少一个仍留在补集的响应，可作为不能保证有限到达的见证。表的构造、真实性和状态访问单独计费。精确商的复合仍满足式~\eqref{pc:eq-exact}，故上述结论沿任意有限商链传递。

近似实现采用逐步交换误差 $\beta_k$、稳定常数 $L$ 和初始误差 $\eta$，则
\begin{equation}
\label{pc:eq-error}
 E_n\le L^n\eta+\sum_{k=0}^{n-1}L^{n-1-k}\beta_k.
\end{equation}
证明是在实际与参考后继之间插入同一输入的实现后继，由三角不等式取得 $E_{k+1}\le LE_k+\beta_k$，再按 $n$ 归纳展开。读回和目标函数的稳定常数依次乘入；离散接受由该误差是否落在已证裕量内决定。

\subsection{被归并的信息、新任务及真实来源}

\begin{theorem}[效用保持、分型恢复与最小细化]
\label{pc:refinement}
给定效用 $U$、分类 $\tau$ 与满射 $q:X\to Q$。$U$ 与 $\tau$ 分别按引理~\ref{pc:factorization} 判定能否读回。若 $U=\bar Uq$ 已成立，并要求追加分类 $\tau$ 且保留旧商全部信息，则
\[
 q^+(x)=(q(x),\tau(x)),\qquad Q^+=\operatorname{im}(q,\tau)
\]
是保留 $q$ 与 $\tau$ 的最粗静态细化。若新增观察改变未来任务，则对增补测试族重新作演化饱和，得到最粗动态充分细化。
\end{theorem}
\begin{certificate}[二坐标回退与新测试]
保存 $q=\operatorname{pr}_1q^+$、$\tau=\operatorname{pr}_2q^+$，以及新观察的来源、版本和实际读取接口。动态证书还保存重新饱和的稳定分区。
\end{certificate}
\begin{proof}
两坐标投影直接恢复旧商和分类。任意其他表示 $p$ 若满足 $q=ap$、$\tau=cp$，则 $h(p(x))=(q(x),\tau(x))$ 良定义，且 $q^+=hp$；$p$ 满射给出唯一性，证明最粗性。动态结论应用定理~\ref{pc:saturation}。新增测试族包含旧族，故新商同像必然旧商同像，存在唯一回退映射；后继交换按代表元代入成立。切换时两种商值来自同一当前输入，阶段内部与切换步骤可按有限前缀归纳合成。
\end{proof}

四点模型 $X=\{0,1\}\times\{b_1,b_2\}$、$q(u,b)=U(u,b)=u$、$\tau(u,b)=b$ 表明，效用可以完全保持而分类消隐。动态模型则取 $a\mapsto c$、$b\mapsto d$、$c\mapsto c$、$d\mapsto d$，并令 $U(a)=U(b)=0$、$U(c)=1$、$U(d)=2$；当前等效用不能决定下一步，加入 $U\Phi$ 后才能区分。因而“无关信息”始终相对于规定任务和时间范围定义；未来行为所需的记忆可以成为当前任务核心的必要组成。

\begin{proposition}[固定接口的信息下限]
\label{pc:information-limit}
若 $q(x_1)=q(x_2)$ 而目标 $T(x_1)\ne T(x_2)$，仅消费 $q$ 的确定性算法不能对两输入都正确；随机算法在两输入上都正确的概率也不能同时严格大于 $1/2$。当目标在度量空间中且两值距离为 $\Delta$ 时，任何单值读回器至少在一输入上误差不小于 $\Delta/2$。
\end{proposition}
\begin{certificate}[同接口异目标见证]
记录两个原输入、相同接口值、不同目标及它们的距离；全部查询和公共参数限定由同一接口给出。
\end{certificate}
\begin{proof}
两输入使确定性执行相同；固定随机种子后仍相同，故随机输出分布相同。两个不同正确值对应互斥事件，其概率和至多为一。单值输出记为 $z$，三角不等式给出 $\Delta\le d(Tx_1,z)+d(z,Tx_2)$，故至少一项达到半距。增加固定接口内的计算次数不改变同像关系，结论继续成立。
\end{proof}

若外部证据 $e:X\to E$ 可合法读取，恢复分类的精确条件改为 $\tau$ 在联合接口 $(q,e)$ 的纤维上恒定。有限域内令 $K_{\max}=\max_y|\tau(q^{-1}(y))|$，全部新增区分由固定长度 $s$ 位提供时必有
\[
 s\ge\lceil\log_2K_{\max}\rceil.
\]
这是同一粗纤维内不同分类必须取得不同码字的计数结论；逐纤维码表可达到该区分界，码表建立、访问和验证仍有独立成本。地址、标签及动态元数据只要能够区分原输入，即计入新增信息；外部证据的构造、存储、访问和验证分别计费。于是，能力增长可以来自有真实来源的合法细化。

\section{率与商的统一：分析微积分、语义余链与同伦见证}
\label{pc:calculus}

“率”描述一个给定承载中的局部变化及沿路径的积累；“商”规定相对于任务可以归并哪些区分，并组织充分表示及其细化。二者的连接需要同一个任务、具体实现映射以及能够同时下降的同伦见证。文献~\cite{S38} 以这一连接把前述法空间、语义场、LHTS与TPH-LHTS的组织原则接到传统求导与积分的严格基础上。

\subsection{两侧承载与显式同伦回缩}

\begin{definition}[分析侧与语义侧]
设 $M$ 为光滑流形，分析侧为 $(\Omega^\bullet(M),d_A)$。设 $(C_\bullet,\partial)$ 为有限有向单纯复形的实链复形，语义侧为
\[
 S^k=\operatorname{Hom}(C_k,\mathbb R),\qquad
 \langle d_Su,c\rangle=\langle u,\partial c\rangle.
\]
顶点表示任务状态，边表示有方向的允许作用，高维链表示已经给出见证的接续关系。几何实现为链映射 $r:C_\bullet\to C_\bullet^{\rm sm}(M;\mathbb R)$，即 $r\partial=\partial r$；生成元的像为有限光滑单形和。
\end{definition}

\begin{theorem}[积分映射的微分交换]
\label{pc:integration}
定义 $\langle\Pi_r\omega,c\rangle=\int_{r(c)}\omega$，则
\[
 d_S\Pi_r=\Pi_r d_A.
\]
对函数 $f$ 与有限路径 $\gamma$，$\langle\Pi_rdf,\gamma\rangle$ 等于 $f$ 与实现后终端零链、起始零链的配对之差；当顶点实现为单点时，即通常的端点函数值之差。
\end{theorem}
\begin{certificate}[逐生成链的Stokes见证]
保存定向、边界矩阵、实现系数和积分定义域；逐生成元核验实现与边界交换。形式在相关紧单形像的邻域光滑。
\end{certificate}
\begin{proof}
对任意相应次数的链 $c$，
\[
 \langle d_S\Pi_r\omega,c\rangle
 =\int_{r(\partial c)}\omega
 =\int_{\partial r(c)}\omega
 =\int_{r(c)}d_A\omega
 =\langle\Pi_rd_A\omega,c\rangle.
\]
第一式是语义微分定义，第二式使用链映射，第三式使用逐单形Stokes公式。有限线性延拓给出算子等式。对一条路径，逐边基本定理使中间零链配对消去，得到两端零链的配对差。
\end{proof}

有限有向图还允许明确逆向插值。把边上光滑且公共端值相容的函数空间记为 $A_G^0$，孤立顶点也保留实值坐标；一形式空间为各边光滑一形式的直和。每条边用 $t\in[0,1]$ 参数化，令 $I^0f$ 读取顶点，$I^1\alpha(e)=\int_0^1\alpha_e$；对顶点数据 $c$ 与边数据 $b$ 定义
\[
 (W^0c)_e(t)=(1-t)c(s(e))+tc(t(e)),\quad
 (W^1b)_e=b(e)\,dt,
\]
\[
 (H\alpha)_e(t)=\int_0^t\alpha_e-t\int_0^1\alpha_e,
 \qquad H|_{A_G^0}=0.
\]
这里 $H\alpha$ 在边两端为零，因而跨公共顶点相容。直接求导和积分得到
\begin{equation}
\label{pc:eq-retraction}
 IW=\operatorname{id},\qquad d_AW=Wd_S,\qquad d_SI=Id_A,
 \qquad \operatorname{id}-WI=d_AH+Hd_A.
\end{equation}
具体地，一形式上 $d_AH\alpha=\alpha-WI\alpha$；零形式上 $Hd_Af=f-WIf$。这证明指定两个复形同伦等价，给出其同调与周期积分的对应。例如 $\alpha=t^2dt$ 时，$I\alpha=1/3$、$H\alpha=(t^3-t)/3$；内部高阶变化由非零见证保存，端点任务与积分同时保持。若新任务要求边内额外细节，则将对应观察加入任务族。

\subsection{棱柱同伦及任务商的共同下降}

\begin{theorem}[棱柱恒等式与商下降]
\label{pc:homotopy-descent}
设两实现 $r_0,r_1$ 由公共面相容的逐单形光滑同伦 $F:[0,1]\times|C|\to M$ 连接。其定向棱柱算子 $P$ 和次数 $-1$ 的积分算子 $K$ 满足
\[
 \partial P+P\partial=r_1-r_0,\qquad
 \langle K\omega,c\rangle=\int_{P(c)}\omega,
\]
\[
 \Pi_{r_1}-\Pi_{r_0}=d_SK+Kd_A.
\]
给定满射链商 $q:C_\bullet\to D_\bullet$，令 $Q^\bullet=\operatorname{Hom}(D_\bullet,\mathbb R)$，并定义共同任务域
\[
 A_T^k=\{\omega:\Pi_{r_0}\omega,\Pi_{r_1}\omega\in\operatorname{im}q_k^*,
                         \ K\omega\in\operatorname{im}q_{k-1}^*\}.
\]
则 $A_T^\bullet$ 对 $d_A$ 封闭，存在唯一 $\Phi_0,\Phi_1,\bar K$，满足
\begin{equation}
\label{pc:eq-main}
\begin{aligned}
 q^*\Phi_i&=\Pi_{r_i},& q^*\bar K&=K,\\
 d_Q\Phi_i&=\Phi_i d_A,&
 \Phi_1-\Phi_0&=d_Q\bar K+\bar Kd_A.
\end{aligned}
\end{equation}
\end{theorem}
\begin{certificate}[棱柱定向与商核条件]
对 $c=[v_0,\ldots,v_k]$ 保存
\[
 P(c)=\sum_{j=0}^k(-1)^j
 F_*[v_0^0,\ldots,v_j^0,v_j^1,\ldots,v_k^1].
\]
逐面配对内部斜面；对 $z\in\ker q$ 分别核验两端积分和棱柱积分为零，以保证映射及见证都与代表元无关。
\end{certificate}
\begin{proof}
交错展开棱柱各单形边界，相邻单形的公共斜面系数相反并消去，侧面合成为 $-P\partial c$，顶面与底面保留 $r_1c-r_0c$。对该等式积分并使用Stokes公式，即得未取商的同伦式。

因为 $q$ 是链映射，$\operatorname{im}q^*$ 对 $d_S$ 稳定。若 $\omega\in A_T$，由积分交换得 $\Pi_i d_A\omega\in\operatorname{im}q^*$；再由
$Kd_A\omega=\Pi_1\omega-\Pi_0\omega-d_SK\omega$ 得 $Kd_A\omega$ 也在该像中，故共同域封闭。余链属于 $\operatorname{im}q^*$ 恰等价于其在 $\ker q$ 上为零，因而可由代表元定义商上算子。满射性使 $q^*$ 单射，保证唯一。将式~\eqref{pc:eq-main} 的后两式复合 $q^*$，化为已证的微分交换及棱柱同伦，再由单射性得到商上等式。
\end{proof}

对闭形式 $d_A\omega=0$，两个实现之差为恰当项，故诱导相同上同调类；在闭链上的积分相同，在一般路径上的差则为 $\langle\bar K\omega,\partial\gamma\rangle$。对零形式观察 $f$，精确观察保持的判据为
\[
 \Phi_1f-\Phi_0f=\bar Kd_Af=0.
\]
对同伦的每个子区间核验这一条件，才得到整条变形中的观察保持。相容细化 $q=rq'$ 同时满足 $\Phi_i'=r^*\Phi_i$、$\bar K'=r^*\bar K$，由复合 $q'^*$ 后的唯一性证明。

非平凡商和非零同伦可以同时存在：取两份单位边，商 $q$ 忘记副本标签；在 $\mathbb R^2$ 中令两份边均由 $(t,0)$ 变形到 $(t,1)$，任务观察为第一坐标。两份副本的积分一致，因而可下降；与此同时
\[
 \langle\bar K(dy\wedge dx),[0,1]\rangle=1,
 \quad \Phi_0(y\,dx)([0,1])=0,
 \quad \Phi_1(y\,dx)([0,1])=1.
\]
任务观察沿变形保持，而见证真实记录了变化。这说明同一任务的不同实现可以有非零内部结构差异。

\subsection{高阶映射与六组联合见证}

在同伦代数层，以悬移空间的约化对称余代数 $S^c(sL)$ 及次数一余导子 $Q_L$ 编码高阶运算，结构闭合为 $Q_L^2=0$。高阶映射 $\mathcal F$ 满足 $Q_B\mathcal F=\mathcal FQ_A$；其一元分量是链映射。若给出次数 $-1$ 的沿映射余导子 $\mathcal H_s$ 及路径
\[
 \partial_s\mathcal F_s=Q_B\mathcal H_s+\mathcal H_sQ_A,
\]
则逐元积分得到 $\mathcal F_1-\mathcal F_0=Q_B\mathcal K+\mathcal KQ_A$。线性任务观察满足 $\tau Q_B=0$、$\tau\mathcal H_s=0$ 时，有 $\partial_s(\tau\mathcal F_s)=0$。复合映射仍与余导子交换，因而有限复合保持高阶关系。非零括号与Jacobiator扇区的高元数据在此独立核验。

\begin{theorem}[率与商的任务保持同伦映射系统]
\label{pc:comparison-main}
在同一任务、源靶对象和共同定义域上，若积分实现、棱柱、商核、观察保持、执行对应与回读、高阶或近似实现的相应条件均成立，则分析微分、语义微分、任务商、同伦变化及有限执行组成相容系统。规定上同调类及其闭链配对、目标可达性、合法策略与终态输出按各自接口保持；新观察通过充分细化加入，近似实现按误差账本和安全裕量接受。
\end{theorem}
\begin{certificate}[六组联合见证]
\begin{enumerate}
 \item 几何实现与边界交换：$r_i\partial=\partial r_i$；
 \item 实现间的棱柱边界与积分同伦；
 \item 两端映射及同伦见证在任务商核上的消去；
 \item 全部规定观察沿同伦的逐值保持；
 \item 合法动作、全部响应、后继及全部接受输出的合法回读；
 \item 使用高阶结构时的逐元关系，以及使用近似实现时的误差界和接受裕量。
\end{enumerate}
各组证书的有效域先取交集，再复合；精确实现不需要额外制造近似误差。
\end{certificate}
\begin{proof}
定理~\ref{pc:integration} 给出微分交换，定理~\ref{pc:homotopy-descent} 使交换式及见证共同下降。闭形式与闭链配对证明同调和周期保持，零形式条件证明任务观察保持。定理~\ref{pc:reach} 的动作响应归纳保持执行、鲁棒可达性与策略回落。高阶映射按余导子交换和路径微分方程复合；近似分支使用式~\eqref{pc:eq-error}。所有证明消费相同任务字段，故可在共同定义域合成。
\end{proof}

LHTS与TPH-LHTS由此取得清楚的位置。棱柱见证说明一次同伦变化的端点差来自何种微分边界；LHTS见证把一次局部扭曲的残差变化归因到指定方向；TPH-LHTS把已认证局部方向组织为递归扩张的候选空间及D结构秩下降；任务商保证全过程保持原任务并允许回读。对残差零形式向量 $\mathbf g$，局部贡献可记为
\[
 \boldsymbol\delta_i^{\partial}=-\langle\bar Kd_A\mathbf g,qv\rangle,
 \qquad \operatorname{span}\mathcal T_k\subseteq\operatorname{span}\mathcal T_{k+1},
 \qquad \rho(D_{t+1})<\rho(D_t).
\]
上述三个式子分别承担局部变化、候选覆盖和实际收敛，不相互替代。它们在联合证书中相接，形成从一致可达、到一致表达、再到递归扩展修复方向的能力链。

\section{局部接续、整体修复与完整成本}
\label{pc:structure-cost}

\subsection{两条独立轴与可计算障碍}

宏观--微观描述同一对象的分辨率关系，局部--整体描述不同区域或接口如何相接。用双索引 $E_{s,i}$ 表示二者：尺度指标 $s$ 配粗化映射，局部指标 $i$ 配覆盖限制和重叠比较，尺度变化与局部限制的交换另作结构条件。这一分工使“增加微观细节”和“补足整体接续”成为能够分别选择、分别证明的操作。

\begin{proposition}[尺度与接续的独立性]
非平凡尺度变化不必产生非平凡局部接续，非平凡接续障碍也不必来自尺度变化。
\end{proposition}
\begin{certificate}[两组反向见证]
尺度例为 $\{0,1\}^2\to\{0,1,2\}$、$(a,b)\mapsto a+b$，覆盖只含整个对象。接续例只有一个尺度，三个局部一维空间的比较为 $g_{12}=g_{23}=1$、$g_{13}=-1$。
\end{certificate}
\begin{proof}
第一例将 $(0,1)$ 和 $(1,0)$ 合并，尺度确实粗化，但单成员覆盖没有跨区域接续义务。第二例没有尺度变化，却有 $g_{23}g_{12}=1\ne-1=g_{13}$ 的接续缺陷。两个例子分别排除两个蕴含，外积构造则使二者同时出现。
\end{proof}

运输残差成为上同调障碍须经历明确构造：先把两条比较路径运至同一纤维，在统一规范中记录有方向坐标 $r_{e,j}$；再指定边、边界与余链 $\omega_j=\sum_e r_{e,j}\varepsilon_e$。图上没有二维面时 $d^1\omega_j=0$；有面时，闭性要求每个面的有向边界残差和为零。顶点势能够修复恰当余链 $\omega=d\eta$；闭而非恰当的余链保留为障碍类。路径上的残差可由前缀累加势恢复，有向环上的非零周期则构成不同义务。

在有限连通图这一具体承载中，$v$ 个顶点、$e$ 条边给出 $G\simeq\bigvee^{e-v+1}S^1$ 及 $H_1(G;\mathbb Z)\cong\mathbb Z^{e-v+1}$。证书先按覆盖顶点数归纳构造含 $v-1$ 条边的生成树，再收缩树边；余下每条边成为一个环，细胞一级边界为零且没有二维边界。因此此类图的同伦正规型可由一级同调秩确定。更一般承载继续保留自身的高阶结构与不变量。

\begin{theorem}[结构障碍、执行修复及取商消去]
\label{pc:obstruction}
设有限维余链复形配正定内积，$d\omega=0$，允许修复集 $\mathcal A\subseteq C^{k-1}$。令
\[
 O_{\rm coh}=\inf_\eta\|\omega-d\eta\|,
 \qquad O_{\rm exec}=\inf_{\eta\in\mathcal A}\|\omega-d\eta\|.
\]
则 $O_{\rm exec}\ge O_{\rm coh}$，且 $O_{\rm coh}=0\iff[\omega]=0$。对微分子复形 $K\subseteq C$ 与商 $p:C\to C/K$，有
\begin{equation}
\label{pc:eq-obstruction-kernel}
 p_*[\omega]=0\quad\Longleftrightarrow\quad
 [\omega]\in\operatorname{im}\bigl(H^k(K)\to H^k(C)\bigr).
\end{equation}
\end{theorem}
\begin{certificate}[核像、投影与原始元]
保存边界矩阵及平方零关系、闭余链、像空间正交基和实际允许修复。商层消去时保存 $p\omega=\bar d\bar h$ 的原始元及其提升 $h$。
\end{certificate}
\begin{proof}
有限维像空间闭合，故到 $\operatorname{im}d$ 的距离由正交投影达到，距离零恰等价于 $\omega\in\operatorname{im}d$。限制优化域使下确界不减。若 $p_*[\omega]=0$，提升商层原始元得 $k=\omega-dh\in K$，且 $dk=0$，于是 $[\omega]$ 来自 $H^k(K)$。反之，若 $\omega=k+dh$ 且 $k\in K$ 闭合，则 $p\omega=\bar d(ph)$，得到商层类为零。
\end{proof}

允许填补的线性化 $B:H\to Z$ 给出同一结论的计算形式：将 $j=j_{\parallel}+j_{\perp}$ 按 $\operatorname{im}B$ 正交分解，则 $\min_h\|Bh-j\|=\|j_{\perp}\|$；可解时，全部解为 $h_0+\ker B$，且 $h_0\perp\ker B$ 为最小范数解。实际执行者还须取得允许的具体 $h$。因此结构可修复、允许算子足够、搜索成功和数值接受是同一修复过程中的不同信息。

若扩张有回缩 $P L\simeq\operatorname{id}$，则 $P_*L_*=\operatorname{id}$，原有非零同调类得以保持。取商消去与保留既有类的回缩修复由此分工明确：取商根据任务归并指定核中的信息，修复增加见证并保持已经承诺的语义。相同残差模长不能决定机制；例如 $I_2$ 与 $\operatorname{diag}(1,-1)$ 具有相同谱范数而迹不同，不能相似，故方向、耦合、时标和算子类别需要各自的观察。

\subsection{修复塔、富结构闭合与超限完成}

设总微分的有限截断为 $D_\varepsilon^{(n)}=D_0+\sum_{r=1}^n\varepsilon^rK_r$，$D_0^2=0$。第 $n$ 阶义务为
\[
 \partial K_n=-O_n,\qquad
 \partial A=[D_0,A],\qquad
 O_n=\sum_{i+j=n,\ i,j\ge1}K_iK_j.
\]
前 $n-1$ 阶闭合时，结合律 $[D^{(n-1)},(D^{(n-1)})^2]=0$ 的第 $n$ 阶系数给出 $\partial O_n=0$，故本阶可解当且仅当 $[O_n]=0$。本阶可解性由命题~\ref{fa:prop-edge}判定，逐阶闭合与形式完成由定理~\ref{ss:total}给出；在有界算子实现中，若 $\sum_{r\ge1}\|K_r\|<\infty$，绝对收敛允许按阶重排双重乘积，从而实际总微分也满足 $D^2=0$。

低阶Jacobiator、完整高阶闭合与工程残差分别登记。一个具有非零二元Jacobiator的完整例子是：$E=\mathbb R^3$，$b(e_1,e_2)=e_2$、$b(e_2,e_3)=e_3$、$b(e_1,e_3)=0$，故 $J_b(e_1,e_2,e_3)=-e_3$。取 $V^{-1}=\widetilde E$、$V^0=E$，定义
\[
 \ell_1\widetilde x=x,\quad \ell_2(x,y)=b(x,y),\quad
 \ell_2(x,\widetilde y)=\widetilde{b(x,y)},\quad
 \ell_3(x,y,z)=\widetilde{J_b(x,y,z)},
\]
其余按分次反对称性延拓，含负次输入的三元项及四元以上运算取零。在约定 $J_{\ell_2}=\ell_1\ell_3$ 的符号下，一、二、三元关系逐项成立；四元关系为三维空间上的交替四线性项而为零；五元项因内部三元输出落在负一次而为零；更高元项至少含一个零运算。于是原二元缺陷非零与总结构闭合相容。

将可缩扇区 $V$ 与任务核心 $T$ 直加，令 $p:T\oplus V\to T$、$i:T\to T\oplus V$，并取 $h(e_j)=\widetilde e_j$，则
\[
 pi=\operatorname{id}_T,\qquad dh+hd=\operatorname{id}-ip.
\]
经 $p$ 因子化的任务保持。这里闭合依赖实际高阶见证；原二元缺陷的大小和修复后剩余方程的大小分别衡量。

\begin{theorem}[相容修复塔的完成与任务保持]
\label{pc:transfinite}
设集合大小的序数塔 $(L_\alpha,j_{\alpha\beta})$ 具有单射结构映射，每个后继修复实际存在并保持已有关系；每项终对象的规定有限元义务均在某阶段取得见证。若极限取相容代数余极限，且任务投影、包含和同伦满足
\[
 p_\alpha i_\alpha=\operatorname{id}_T,\quad
 d_\alpha h_\alpha+h_\alpha d_\alpha=\operatorname{id}-i_\alpha p_\alpha,
\]
并与塔映射交换，则终对象保持这些关系及任务回读。拓扑完成另要求相容的完备分离滤过和运算连续性。
\end{theorem}
\begin{certificate}[初始、后继、极限及义务覆盖]
每项义务记录出现阶段、填补阶段与后继运输；代数极限保存有限输入的共同阶段。完成保存 $L\simeq\varprojlim_pL/F^pL$、$\bigcap_pF^pL=0$ 及
$\ell_n(F^{p_1},\ldots,F^{p_n})\subseteq F^{p_1+\cdots+p_n}$。
\end{certificate}
\begin{proof}
作超限归纳。初始由输入证书成立，后继由实际填补及保持式成立。极限阶段的每组有限输入来自某一共同早期阶段，运算和有限项关系在那里求值，再映入余极限；代表元更换可在更晚共同阶段比较，故定义良好。每项义务由覆盖证书在某阶段闭合，并被输送到终对象。任务映射同样按代表元诱导，逐式代入得到 $pi=\operatorname{id}$ 及同伦式。完备情形在每个截断商验证，完备性给出存在，分离性给出唯一与等式成立。规定有限运行再按路径长度归纳保持。
\end{proof}

超限阶段、运算元数和数值精度是三个独立参数。对集合大小证书域 $S$ 的单调膨胀闭包，反复加入新证书必在 $|S|^+$ 之前稳定；若所有规则只有有限前提，$\omega$ 阶并集已经闭合，因为每个有限前提集都在某一有限阶段共同出现。若有真正无限前提，则可能出现 $\omega+1$ 阶才产生的新义务。以上是完成对象及关系保持的存在性组织；某个任务的有限运行由其具体秩下降、阈值及允许操作另行确定。

\subsection{母结构增长、任务核心与真实工作量}

令任意有限维基余链复形为 $(B,d_B)$，并构造
\[
 C_N=B\otimes\Lambda(\xi_1,\ldots,\xi_N),\qquad
 D_N=d_B\otimes1,\qquad |\xi_i|=1,
\]
则按外代数单项式分块，核和像分别为基复形核像的直和，故
\[
 H^\bullet(C_N)\cong H^\bullet(B)\otimes\Lambda_N,
 \qquad \dim H^\bullet(C_N)=2^N\dim H^\bullet(B).
\]
每增加一生成元，$\Lambda_{N+1}=\Lambda_N\oplus\Lambda_N\xi_{N+1}$，给出维数翻倍的归纳证书。增广 $p_N=\operatorname{id}\otimes\epsilon$ 保留常数块，核为 $B\otimes\Lambda_N^{>0}$。全部任务测试经 $p_N$ 因子化时，附加循环可从任务表示中归并，核心仍完整。这一例子证明母上同调的增长与有效任务核心的大小可以独立变化。

特别取 $B=\mathbb R$ 集中于零次、$d_B=0$、$N=1$。增广
$p:\mathbb R\oplus\mathbb R\xi\to\mathbb R$ 保持全部经 $p$ 因子化的任务，但
$H^1(\mathbb R\oplus\mathbb R\xi)=\mathbb R$ 而 $H^1(\mathbb R)=0$，故 $p$ 不是链同伦等价。这给出任务充分性与整体同伦等价各有独立见证的具体实例。

状态层可相应构造 $X_N=S\times Z_N$、$q_N(s,z)=s$；若动作、响应、目标和测试只依赖 $s$，后继第一坐标为 $T(s,a,\xi)$，则 $q_N$ 是精确任务商。链层增广、状态层任务商与任务测试的图实现分别给出，再以交换式连接。严格算子层取分次交换子及微分 $\partial a=Da-(-1)^{|a|}aD$；$D^2=0$ 给出 $\partial^2=0$，结合律给出分次Jacobi，保持商核的算子可下降到商上。主丛提供其等变承载，不能代替这些逐层保持条件。

\begin{theorem}[统一接口的总成本及复用收益]
\label{pc:cost}
给定已认证任务商和核心求解器，一次事件按顺序执行核心提取、修复、核心求解、合法回落、输出与验证，成本分别为 $c_q,c_h,c_s,c_r,c_o,c_v$。初始构造认证成本为 $B$，运行 $K$ 次的总工作量为
\[
 W(K)=B+\sum_{t=1}^K
 (c_{q,t}+c_{h,t}+c_{s,t}+c_{r,t}+c_{o,t}+c_{v,t}).
\]
对同一输入权限、成功条件和输出义务，若基准每次成本为 $a$，新接口每次为 $c$，则严格节省成立当且仅当 $K(a-c)>B$。
\end{theorem}
\begin{certificate}[不重复计费与有效区间]
固定成本原语，把每个实际原语恰计入一个阶段，记录证书来源、任务、规则和有效期；变更引起的细化与重新认证列为新成本。并行墙钟时间另按调度计算。
\end{certificate}
\begin{proof}
串行阶段的工作量相加得到首式。常成本情形 $W(K)=B+Kc$，与 $Ka$ 比较并移项即得等价式；平均成本为 $c+B/K$。使用上下界时，相同严格不等式给出充分收益判据。任务与规则改变后按有效区间分段求和即可。
\end{proof}

小商的存在与取得小商值的成本分别计量。奇偶任务 $q_N(x)=x_1\oplus\cdots\oplus x_N$ 只有两个商值，但任何逐位查询的确定性精确算法必须读取全部 $N$ 位：少读一位时翻转该位保持所有已读值，却改变答案。已知认证商值且每次提供一个改动位的旧、新值时，维护式 $q_{\rm new}=q_{\rm old}\oplus u\oplus v$ 才有常数操作成本。

同样，固定 $d$ 时，$X_N=\mathbb F_p^d\times\mathbb F_p^N$ 上只求 $F(c,z)=\sum_{i=1}^dc_i$，恰需 $d$ 次核心坐标查询，相对扩展规模 $N$ 为常数；每个核心纤维仍可含 $p^N$ 个不同附加分型。若任务要求完整恢复 $z$，其附加区分至少为 $\lceil N\log_2p\rceil$ 位；每次至多输出 $b$ 个坐标时，完整输出至少需 $\lceil N/b\rceil$ 次原语。$n$ 个变量到总阶数 $k$ 的完整Taylor表有 $\binom{n+k}{n}$ 个系数槽，物化这些槽承担相应输出成本。

固定字长、固定核心字段数 $m$ 的查询模型中，已给总的核心求解器便给出至多 $m$ 次读取的上界；若两个合法输入的正确输出集合不相交，零查询无法区分，确定性与逐输入成功率大于 $1/2$ 的量子查询都至少需一次。因而两者相对非干扰扩展参数均为 $\Theta(1)$，这一结论的参数是\emph{固定核心接口}；提取、认证、输出和实际算术工作量继续在完整账本中出现。

\section{哲学收束：数学有效性、表达演化与可审计的实践闭环}
\label{pc:philosophy}

\subsection{对象、意义与同一性：以任务保持连接不同承载}

G框架的哲学起点，是把对象、规律、变化、观察及其相容条件放在同一个可追溯体系中。对象可以具有不同的集合编码、几何实现、语义表达和程序实现；判断它们是否承担同一个任务，要检查观察、合法作用、响应、目标与结果回读。任务保持由具体映射实现，同伦见证说明允许改变的内部结构如何沿一条有类型路径连接。因而意义具有明确的关系性：它存在于对象能够支持的区分、作用及回读关系中，并由共同任务的证书固定。

这一立场同时容纳稳定性与变化。性质塔保留已经声明的观察，新的性质使表示更细；同伦修复增加接续见证，原任务经回缩继续保持；取商则把相对于当前任务等价的表示归并。三个操作的结果不同，恰由它们各自的映射确定。将一个对象从一种表达改写为另一种表达时，能够继续使用的结论来自保持证书，新增的能力来自新接口和新结构的实际构造。

因此，框架中的“同一”始终带有明确限定：状态同一性由测试等价给出，映射同伦由边界填补给出，整体同伦等价需要互逆映射及见证，数值近似由误差界给出。它们形成可相接的层次。这样的同一性使跨尺度、跨语言、跨应用的迁移可以具体比较，也让各层的变化保留自己的信息量。

\subsection{数学为什么有效：从共同不变量到可检验回读}

数学有效性的核心，在于现实任务的某些关系可以由数学对象保持，而数学操作又与真实允许操作相容。定理~\ref{pc:wigner} 给出了这一点的有限执行形式；率--商主定理进一步加入积分、同伦、细化和高阶相容。数学侧计算的量能够回到原任务，依靠的是观察因子化、后继交换、稳定性及合法回读。这使“有效”成为一组可以逐项检查的条件。

共同正规型的价值，在于为原本异质的表示建立共同可读的核心。物理观测、语义对象、代数运算和代码执行各有自己的来源和类型；只要在指定任务上保持相同关系，它们就能够共享同一条证明链。源证据使模型指向具体对象，数学结构使关系能够推演，算法使构造能够执行，回读检验使执行重新接受对象约束。四者的连接形成解释与实践之间的闭环。

对涌现与复杂系统，这一闭环允许选择真正需要的宏观输出。宏观性质若已由有效核心、不变量和边界证书决定，求解可以集中于该核心；若任务后来要求更多机制、轨迹或微观分类，则新的输出和区分进入表示及成本。数学的压缩能力由任务充分性支撑，数学的解释能力由保留关系支撑，两者在同一个表示体系中可以共同增长。

\subsection{率与商的互补：变化计算与表示组织共同决定能力}

率计算给定承载中的变化，商决定哪些差异在当前任务下需要保留。连续分析通过导数、积分、极限和误差估计处理变化；语义组织通过生成元、关系、任务测试和高阶填补处理承载之间的连接。前者提供局部变化的可积累结构，后者提供这些变化可以比较、归并和传递的条件。式~\eqref{pc:eq-main} 把两个方向接成同一组算子恒等式。

“离散组织连续”在这里具有具体含义：离散生成元、边界和接口规定连续纤维及其接续；连续积分又把纤维内部变化送入离散语义链。有限图中的插值与同伦回缩表明，两个方向可以有显式往返；更一般情形则由共同定义域和商核条件决定可传递的内容。连续性和结构相容性分别贡献自己的严密性。

宏观--微观与局部--整体的双轴还指出两类改善方向。分辨率细化增加可识别的自由度、精度和耦合层次；整体修复补足局部呈示之间的相容关系，允许在任务核心上组织求解。需要更多细节的任务沿第一轴细化，需要消除接续障碍的任务沿第二轴修复，综合任务可以同时进行。增加相容阶数与减少有效状态数分别有自己的证明，不能仅由表面维数互相推断。

\begin{corollary}[有来源的表示细化产生严格能力增量]
\label{pc:capability}
在输入 $X=\mathbb R\times\{0,1\}$ 上，若固定接口只给 $q(t,b)=t$，则目标 $T(t,b)=b$ 无法精确求解。若允许从同一真实输入读取并认证分支位，则细化接口 $q'(t,b)=(t,b)$ 能精确求解，并保持旧接口全部可解任务。
\end{corollary}
\begin{certificate}[来源、回退及目标回读]
记录分支位的有效读取、类型及认证，粗化映射为 $r(t,b)=t$，满足 $rq'=q$；目标回读为第二坐标投影。
\end{certificate}
\begin{proof}
固定 $t$ 的两个输入是命题~\ref{pc:information-limit} 的同像异目标对。新接口直接给出分支位，因而恢复目标；任意旧接口算法可以先施加 $r$ 再照常执行，所以原能力全部保留，而此目标新增可解，得到严格包含。
\end{proof}

这给出了表达演化的精确认识论：既定表示的边界可以通过合法增加证据、区分和结构被扩展。扩展的贡献由新输入与新规则解释，旧结论则经回退和保持式延续。由此，研究过程可以围绕“哪个区分缺失、哪个接续未闭合、哪项见证尚未取得”持续前进。

\subsection{障碍的知识地位：从误差数值到可修复结构}

障碍提供的是有结构的信息。观测残差说明实现与任务之间存在差异；经过运输、复形和闭性核验的障碍类，进一步说明哪些差异不能由当前允许的恰当变化消除。剩余障碍的方向、周期、边界和依赖关系决定应当修改数据、拓展算子、细化任务商还是补足高阶关系。这样的分类使失败能够回到构造过程，形成新的研究对象。

局部成功进入整体结论，需要把局部见证真正接起来。棱柱恒等式说明变化的来源，LHTS记录局部方向的贡献，TPH-LHTS组织候选方向和递归秩，任务商保持目标与回读。每一环提供下一环所需的输入，其合取构成整体修复。形式完成、滤过完成和有限执行继续使用各自的见证，从而把理论存在、可计算实现和实际完成放入同一个分层体系。

证书的作用也由此明确：它记录一项结论怎样成立，以及在什么输入、规则和定义域内继续成立。构造定理把可证明的系统规则从先验断言转为生成式结论；原子化审计把复杂表达分解为可以检查的依赖；未决项保留其位置，已完成项能够复用。知识增长表现为证书图的有来源扩张与修复，既有真值不因新的表达需求而被无依据改写。

\subsection{有效性、成本与治理：把任务承诺落实到全路径}

效用保持确定某个输出是否继续有效，分类保持确定某些机制、来源和规则能否继续辨认。两者分别因子化，使当前任务中的压缩与未来任务中的责任可以明确安排。若某种分类属于必须保留的治理要求，它就进入测试族和读回接口；若其暂不属于当前输出，也可以由独立证据层保留，以便后续任务据此细化。这个安排直接使用联合接口的充分性和信息下界。

\begin{remark}[能力、用途与来源证据的各自职责]
同一种构造能力可以服务于不同目标；保护量、权限和发布条件作为任务字段参与保持与认证。DOI、时间戳和版本记录承担来源、优先权及复现定位，具体命题的成立仍由证明承担，实验结论由相应数据与验收承担。各类证据在同一证书链中协作，并保留自身的证明职责。\cite{S14,S16,S37,S38}
\end{remark}

目标、效用、合法作用及允许损失由任务给定；数学证明核验这些承诺，任务变更触发新的表示与认证。路径治理由此落实为逐步记录合法性、保持量、剩余障碍、回读与成本，分别保留数值收益、分类信息和来源证据。

效率评价覆盖源数据读取、商构造、修复、核心计算、验证、输出和重新认证。充分商与不变量复用的收益由定理~\ref{pc:cost}核算；结构能力的增加和实际性能的提高分别给出证书，并在同一任务上比较。

全文以任务充分商和率--商比较完成哲学收束：统一性由可组合的保持关系成立，开放性由有来源的细化与修复实现；数学解释与实践经证书、成本和合法回读相接。


\clearpage
\appendix
\section{统一术语与符号索引}
\label{app:notation}
下表汇集跨章节反复出现的术语；具体数学对象的域、次数和适用条件仍以首次定义为准。
\small
\begin{longtable}{P{0.30\textwidth}P{0.64\textwidth}}
\toprule 记号或名称 & 在本文中的职责\\\midrule
\endfirsthead
\toprule 记号或名称 & 在本文中的职责（续）\\\midrule
\endhead
\midrule\multicolumn{2}{r}{续下页}\\\endfoot
\bottomrule\endlastfoot
PL--PI & 泛逻辑与泛迭代；分别组织判定/偏好和异构演化及规则更新。\\
GZ-GMS & 高政广义数学结构；共同登记对象、规律、流、性质层、联络与曲率。\\
GZ-LS / LT / LOC / LCurv / TLC & 法空间、法变换、法联络、法曲率及三层联络。\\
PFB-GNLA / G代数 & 主纤维丛上的广义非交换代数及其严格、同伦、算子泛函和元数学版本。\\
GZ-OHU & 算子--同伦普适性（Operator--Homotopy Universality）；以逐阶段附加、保持及完成见证承载修复。\\
D结构 / DERI & 决策路径和递归基准更新；DERI承担声明环境与历史下的基准拟合。\\
GRL / KAT / MDQ & 法空间路径学习；含测试的Kleene代数；可独立检查和组合的微分动力量子。\\
HACA / PACER & 分层代数认知架构；词包对齐的压缩--展开推理闭环。\\
LBOPB & 生命总算子主纤维丛；七算子域和六观察框架在共同生命对象上协变。\\
PTOPB / ETOPB & 物理/工程总算子丛；联合承载状态、构造、控制、制造、验证和运维。\\
ORL & 开放系统中的观察、行动、残差及结构修复学习；接口与开放等级显式给定。\\
$\mathsf L_1,\mathsf L_2,\mathsf L_{\ge3}$ & 离散索引、连续基底及高层修复纤维；与L测度的分辨率编号区分。\\
LHTS / TPH & 格点同伦扭曲评分与轨迹函数参数同伦；分别提供局部运动和外层选择。\\
Jacobiator-gap / $\Phi$ & 在声明类型中计算的闭合缺口及其势函数；具体分量须与任务绑定。\\
$C_\bullet,\partial;\ C^\bullet,d$ & 链复形与余链复形；平方零使同调/上同调的核像商良定义。\\
$\ell_n,Q$ & 高阶同伦操作及编码它们的余导子；$Q^2=0$组织全部元数的闭合关系。\\
$q:X\to Q_{\mathcal T}$ & 任务表示或充分商；纤维恒定性、后继交换和合法回读决定保持范围。\\
$q^*$ & 链商诱导的余链拉回；满射链商对应单射余链映射。\\
$\Pi_r,\Phi_i,K,\bar K$ & 几何积分映射、商上比较映射、棱柱积分及其下降后的同伦见证。\\
$\mathcal C,\mathsf{Open},V$ & 证书、未解除前提集与验证器；结论绑定输入、规则、任务及有效域。\\
$c_q,c_h,c_s,c_r,c_o,c_v$ & 商提取/维护、修复、核心求解、回落、输出和验证的实际成本。\\
\end{longtable}
\normalsize
\clearpage
\section{三十八篇源文稿的覆盖与整合索引}
\label{app:coverage}
编号按附件中的出现顺序排列；主题采用简名。正文按概念和证明依赖合并重复材料，下表标出各稿核心内容的主要归属；精确文件名与原文行域见参考文献。
\small
\begin{longtable}{P{0.055\textwidth}P{0.235\textwidth}P{0.16\textwidth}P{0.47\textwidth}}
\toprule 源稿 & 主题简名 & 本文位置 & 核心保留内容\\\midrule
\endfirsthead
\toprule 源稿 & 主题简名 & 本文位置 & 核心保留内容（续）\\\midrule
\endhead
\midrule\multicolumn{4}{r}{续下页}\\\endfoot
\bottomrule\endlastfoot
\cite{S01} & 纯数学：PL--PI、法空间与G代数 & \ref{fa:foundation}、\ref{fa:algebra}节 & GZ命名、三层联络、四种G代数、CH分层、OHU与严格结构\\
\cite{S02} & 应用数学I：物理、GRL、量子与材料 & \ref{fa:applications}节 & 法空间作用量、策略路径、量子接口、超导与超流研究任务\\
\cite{S03} & 应用数学II：生命总算子丛 & \ref{fa:applications}节 & 七幺半群、六观察框架、宏微协变、逆向设计与虚拟人体\\
\cite{S04} & 应用数学III：HACA/PACER & \ref{fa:semantics}节 & 性质层、词包KAT、双GRL、MDQ、路径证书与三层完备性\\
\cite{S05} & Jacobi缺口与D结构等价 & \ref{fa:algebra}节 & DERI、路径同序商、D-Decision及D-Depth\\
\cite{S06} & 同伦、间隙与最小性 & \ref{fa:foundation}、\ref{fa:algebra}节 & 任务分辨率、最小性所依赖的基准与迭代深度\\
\cite{S07} & 高阶同伦修复塔 & \ref{fa:algebra}、\ref{pc:structure-cost}节 & 高元关系、障碍类、修复深度、广度与完成\\
\cite{S08} & 边缘算子与变分塔 & \ref{fa:algebra}节 & 新算子来源、BV主方程与高阶闭合\\
\cite{S09} & 同伦代数、公理与语义工作流 & \ref{fa:foundation}、\ref{fa:algebra}节 & 构造式定理、封闭证书、性质变换及修复程序\\
\cite{S10} & 同伦同一性与单子集合论 & \ref{fa:foundation}、\ref{pc:quotient}节 & 性质塔、生成式对象、单子实例与见证重建\\
\cite{S11} & 语义函子场与规范化接口 & \ref{fa:semantics}节 & 部分规范化、标准族、语义与风格的回读关系\\
\cite{S12} & 正则规范化、验证及鲁棒性 & \ref{fa:semantics}节 & 候选收缩、MDQ运行、双指纹、重复验证与跨实现一致性\\
\cite{S13} & PTOPB与逆向设计 & \ref{fa:openengineering}节 & 相内搜索、包络迁移、联合控制、可达性与代价下界\\
\cite{S14} & PTOPB/ETOPB算子分解 & \ref{fa:openengineering}、\ref{pc:philosophy}节 & 物理/工程与内生/外生双分类、扇区、工程生命周期及用途\\
\cite{S15} & 开放完备等级与RTSC治理 & \ref{fa:openengineering}节 & 扰动类、修复预算、未知事件与固定观察边界\\
\cite{S16} & 幂域纤维扩展与ORL & \ref{fa:openengineering}节 & 离散规范场、宏态纤维、路径权重、开放学习与独立证据\\
\cite{S17} & 缺口动力与统计力学下降 & \ref{ss:dynamics}节 & 门控、积分触发、阻尼、预算、组织同伦与实际部署\\
\cite{S18} & 可修复障碍与广义分形 & \ref{ss:repair}节 & 修复对象域、尺度递归子类、超限单子、卡丘胚与核商\\
\cite{S19} & 主联络、路径积分与同伦扭曲 & \ref{ss:semantic}节 & 水平提升、终点商、支撑类及正规形三重编码\\
\cite{S20} & 语义规范场与自然语言运行时 & \ref{ss:semantic}节 & 规范语义、逻辑占位、文本回读、插件及白盒主权\\
\cite{S21} & 离散修复塔与量子几何 & \ref{ss:repair}节 & 胞腔填补、同调秩、拉普拉斯谱及量子实现接口\\
\cite{S22} & 扩展签发塔 & \ref{ss:semantic}节 & 证书扩充、元素同调表示、可数塔编码与访问成本\\
\cite{S23} & 截面孔洞与条件收敛 & \ref{ss:repair}节 & 全局截面障碍、势函数、下降与极限条件\\
\cite{S24} & Jacobi--Bianchi离散修复 & \ref{ss:repair}节 & 缺陷登记、终止排名、预算嵌套与终端正规形\\
\cite{S25} & 生命/语义开放修复引擎 & \ref{ss:open}节 & V2/V3接口、统计选类、领域迁移与有限交付\\
\cite{S26} & 离散--连续--修复三级结构 & \ref{ss:repair}节 & 可数索引、连续基底、修复纤维、视域与双重收敛\\
\cite{S27} & 宏观--微观总同调 & \ref{ss:repair}节 & 水平/垂直耦合、总微分、时间结构、BA交替与嵌入修复\\
\cite{S28} & 观测、推断与全局截面 & \ref{ss:open}节 & 观察因子化、不可辨识、边缘算子与合法证据扩张\\
\cite{S29} & Forcing、性质态射与修复 & \ref{ss:open}节 & 扩展承载、高阶闭合、旧观察保真及有限交付\\
\cite{S30} & 真理锚点与语义障碍 & \ref{ss:semantic}节 & 锚点、标准变更、有限语义链保持与反身审计\\
\cite{S31} & 拒绝签发与概念降级 & \ref{ss:semantic}节 & 未闭合项定位、条件化命题、语义附加与修复入口\\
\cite{S32} & 统计格点与PDE策略簇 & \ref{ss:dynamics}节 & 方法空间、Gibbs分布、离散切换及近似量子软件接口\\
\cite{S33} & 同伦命中与PDE范式替代 & \ref{ss:hitting}节 & 首次命中、条件概率、幂化关系、高阶正交与厚通道\\
\cite{S34} & LHTS及延迟信用 & \ref{ss:hitting}节 & 分量评分、持久信用、雅可比反拉回及覆盖义务\\
\cite{S35} & TPH--LHTS轨迹参数同伦 & \ref{ss:hitting}节 & 概率路径、递归轨迹族、预测裕量、回填与D结构下降\\
\cite{S36} & PDE的L测度与隧穿等价 & \ref{ss:pde}节 & 测度微分、弱残差、表达/求解/解释迁移及工程归因\\
\cite{S37} & 基础障碍谱系 & \ref{pc:spectrum}、\ref{pc:philosophy}节 & 完整希尔伯特索引、跨学科桥接、原子推理与数学有效性\\
\cite{S38} & 率与商 & \ref{pc:quotient}、\ref{pc:calculus}、\ref{pc:structure-cost}、\ref{pc:philosophy}节 & 最粗充分商、六组见证、双轴、动态记忆与证据细化、超限、完整成本及哲学闭环\\
\end{longtable}
\normalsize
\clearpage
\phantomsection
\addcontentsline{toc}{section}{参考文献与原稿定位}
\begin{thebibliography}{99}
\setlength{\itemsep}{0.65em}
\small
\bibitem{S01}
高政（GaoZheng）. Meta-Mathematical Theory based on Pan-Logic Analysis and Pan-Iterative Analysis (GaoZheng G-Framework) and the Principal-Bundle-Based Generalized Noncommutative Lie Algebra (GaoZheng G-Algebra): An Integrated Construction. 2025.
\newblock Zenodo，引用版本 v1.0；本地源文件标注 Version 1.1, 2025。
\newblock DOI：\url{https://doi.org/10.5281/zenodo.17651584}。
\newblock 本地源文稿：\path{docs/project_docs/Pub_GFramework_PureMath.tex}，第1--15562行。
\bibitem{S02}
高政（GaoZheng）. GaoZheng G-Framework and GaoZheng G-Algebra: Applied Mathematics Volume I -- Law-Space Geometry, GRL, Quantum Computing, and Superconductivity. 2025.
\newblock Zenodo，引用版本 v1.2；本地源文件标注 Version 1.3, 2025。
\newblock DOI：\url{https://doi.org/10.5281/zenodo.17686133}。
\newblock 本地源文稿：\path{docs/project_docs/Pub_GFramework_AppMath_Vol1.tex}，第1--7532行。
\bibitem{S03}
高政（GaoZheng）. GaoZheng G-Framework and GaoZheng G-Algebra: Applied Mathematics Volume II -- LBOPB, Life-Science Monoids, and Generative Precision Medicine. 2025.
\newblock Zenodo，引用版本 v1.1；本地源文件标注 Version 1.1, 2025。
\newblock DOI：\url{https://doi.org/10.5281/zenodo.17744672}。
\newblock 本地源文稿：\path{docs/project_docs/Pub_GFramework_AppMath_Vol2.tex}，第1--5421行。
\bibitem{S04}
高政（GaoZheng）. GaoZheng G-Framework and GaoZheng G-Algebra: Applied Mathematics Volume III -- HACA, PACER, and Certificate-Based AI. 2025.
\newblock Zenodo，引用版本 v1.0；本地源文件标注 Version 1.0, 2025。
\newblock DOI：\url{https://doi.org/10.5281/zenodo.17762295}。
\newblock 本地源文稿：\path{docs/project_docs/Pub_GFramework_AppMath_Vol3.tex}，第1--6924行。
\bibitem{S05}
高政（GaoZheng）. 雅可比间隙与 D 结构在语义层面的严格等价性（工作笔记）. 2025-12-15.
\newblock 收录于《中文工作笔记》v5.2；文集 DOI：\url{https://doi.org/10.5281/zenodo.22739417}。
\newblock 本地源文稿：\path{docs/project_docs/notebook_tex_v5.2/20251215_jacobi_gap_D_structure_equivalence.tex}，第1--2099行。
\bibitem{S06}
高政（GaoZheng）. 同伦、雅可比间隙与离散统本体：结构化笔记. 2025-12-20.
\newblock 收录于《中文工作笔记》v5.2；文集 DOI：\url{https://doi.org/10.5281/zenodo.22739417}。
\newblock 本地源文稿：\path{docs/project_docs/notebook_tex_v5.2/20251220_homotopy_minimality_dialogue.tex}，第1--3209行。
\bibitem{S07}
高政（GaoZheng）. 在 G 框架 中：更高阶 $L_\infty$ 算子作为同伦修复塔 （纯形式系统笔记）. 2025-12-28.
\newblock 收录于《中文工作笔记》v5.2；文集 DOI：\url{https://doi.org/10.5281/zenodo.22739417}。
\newblock 本地源文稿：\path{docs/project_docs/notebook_tex_v5.2/20251228_gframework_linfty_homotopy_repair_formal_system.tex}，第1--2497行。
\bibitem{S08}
高政（GaoZheng）. 边缘算子封闭性、比安基基准与变分生成的同伦修复塔 （工作笔记：形式系统版）. 2026-01-11.
\newblock 收录于《中文工作笔记》v5.2；文集 DOI：\url{https://doi.org/10.5281/zenodo.22739417}。
\newblock 本地源文稿：\path{docs/project_docs/notebook_tex_v5.2/20260111_gframework_edge_operator_homotopy_closure_variational_tower_formal_system.tex}，第1--773行。
\bibitem{S09}
高政（GaoZheng）. 性质代数与同伦代数的封闭证书、障碍类与拓扑修复 （工作笔记：形式系统版）. 2026-01-12.
\newblock 收录于《中文工作笔记》v5.2；文集 DOI：\url{https://doi.org/10.5281/zenodo.22739417}。
\newblock 本地源文稿：\path{docs/project_docs/notebook_tex_v5.2/20260112_gframework_homotopy_algebra_axiom_semantics_workflow_formal_system.tex}，第1--983行。
\bibitem{S10}
高政（GaoZheng）. 同伦同一性、性质集与单子集合论：离散统连续统分层的证书化奠基 （工作笔记：形式系统版）. 2026-01-12.
\newblock 收录于《中文工作笔记》v5.2；文集 DOI：\url{https://doi.org/10.5281/zenodo.22739417}。
\newblock 本地源文稿：\path{docs/project_docs/notebook_tex_v5.2/20260112_gframework_homotopy_identity_monadic_set_theory_foundation_formal_system.tex}，第1--1132行。
\bibitem{S11}
高政（GaoZheng）. 语义函子场与语义规范场：词包判定--规范化接口 （工作笔记：形式系统版）. 2026-01-13.
\newblock 收录于《中文工作笔记》v5.2；文集 DOI：\url{https://doi.org/10.5281/zenodo.22739417}。
\newblock 本地源文稿：\path{docs/project_docs/notebook_tex_v5.2/20260113_gframework_semantic_functor_field_gauge_field_normalization_interface_formal_system.tex}，第1--1298行。
\bibitem{S12}
高政（GaoZheng）. 语义函子：正则生成、裁决与规范化接口 （工作笔记：形式系统版）. 2026-01-14.
\newblock 收录于《中文工作笔记》v5.2；文集 DOI：\url{https://doi.org/10.5281/zenodo.22739417}。
\newblock 本地源文稿：\path{docs/project_docs/notebook_tex_v5.2/20260114_gframework_semantic_functor_regex_normalization_validation_robustness_formal_system.tex}，第1--2693行。
\bibitem{S13}
高政（GaoZheng）. Applied Math Vol II：从 LBOPB 到 PTOPB 的可迁移模板 （工作笔记：形式系统版）. 2026-01-16.
\newblock 收录于《中文工作笔记》v5.2；文集 DOI：\url{https://doi.org/10.5281/zenodo.22739417}。
\newblock 本地源文稿：\path{docs/project_docs/notebook_tex_v5.2/20260116_gframework_appmath_vol2_ptopb_inverse_design_formal_system.tex}，第1--8384行。
\bibitem{S14}
高政（GaoZheng）. PTOPB/ETOPB 的算子二分：外界实时干扰与非交换构造 （工作笔记：形式系统版）. 2026-01-18.
\newblock 收录于《中文工作笔记》v5.2；文集 DOI：\url{https://doi.org/10.5281/zenodo.22739417}。
\newblock 本地源文稿：\path{docs/project_docs/notebook_tex_v5.2/20260118_gframework_ptopb_etopb_operator_decomposition_formal_system.tex}，第1--13293行。
\bibitem{S15}
高政（GaoZheng）. 开放系统博弈的完备性：等级、同伦修复塔与工业化室温超导治理 （工作笔记：形式系统版）. 2026-01-25.
\newblock 收录于《中文工作笔记》v5.2；文集 DOI：\url{https://doi.org/10.5281/zenodo.22739417}。
\newblock 本地源文稿：\path{docs/project_docs/notebook_tex_v5.2/20260125_gframework_open_system_completeness_levels_rtsc_governance_formal_system.tex}，第1--1286行。
\bibitem{S16}
高政（GaoZheng）. 幂域纤维扩张与 ORL 规范场：从相空间扩张到同调修复与 OHU 正交子丛 （工作笔记：形式系统版）. 2026-01-30.
\newblock 收录于《中文工作笔记》v5.2；文集 DOI：\url{https://doi.org/10.5281/zenodo.22739417}。
\newblock 本地源文稿：\path{docs/project_docs/notebook_tex_v5.2/20260130_gframework_powerdomain_fiber_extension_orl_gauge_field_formal_system.tex}，第1--9458行。
\bibitem{S17}
高政（GaoZheng）. 基于雅可比间隙动力驱动的同伦扭曲与统计力学最小势能路径收敛. 2026-02-17.
\newblock 收录于《中文工作笔记》v5.2；文集 DOI：\url{https://doi.org/10.5281/zenodo.22739417}。
\newblock 本地源文稿：\path{docs/project_docs/notebook_tex_v5.2/20260217_gframework_jacgap_homotopy_distortion_statmech_min_potential_path_convergence.tex}，第1--6173行。
\bibitem{S18}
高政（GaoZheng）. 可修复障碍同伦类与广义分形：全集范畴与严格边界 （工作笔记：形式系统版）. 2026-02-21.
\newblock 收录于《中文工作笔记》v5.2；文集 DOI：\url{https://doi.org/10.5281/zenodo.22739417}。
\newblock 本地源文稿：\path{docs/project_docs/notebook_tex_v5.2/20260221_gframework_repairable_obstruction_homotopy_class_vs_generalized_fractal_formal_system.tex}，第1--6751行。
\bibitem{S19}
高政（GaoZheng）. 主丛联络的纤维丛空间、路径积分支撑类与同伦扭曲正规形的集合等价 （工作笔记：形式系统版）. 2026-03-09.
\newblock 收录于《中文工作笔记》v5.2；文集 DOI：\url{https://doi.org/10.5281/zenodo.22739417}。
\newblock 本地源文稿：\path{docs/project_docs/notebook_tex_v5.2/20260309_gframework_principal_connection_fiber_space_path_integral_homotopy_twist_equivalence_formal_system.tex}，第1--5791行。
\bibitem{S20}
高政（GaoZheng）. G 框架正则语义规范场到自然语言语义逻辑占位规范场的 高阶同伦扭曲双射、LLM 运行时降级与逻辑引擎 （工作笔记：形式系统版）. 2026-03-10.
\newblock 收录于《中文工作笔记》v5.2；文集 DOI：\url{https://doi.org/10.5281/zenodo.22739417}。
\newblock 本地源文稿：\path{docs/project_docs/notebook_tex_v5.2/20260310_gframework_semantic_gauge_field_natural_language_logic_placeholder_runtime_formal_system.tex}，第1--16233行。
\bibitem{S21}
高政（GaoZheng）. 离散统的修复塔与量子--几何接口 （从 $H_1$ 收缩到扇区减少与离散 Hodge 谱隙的可审计证明稿）. 2026-03-13.
\newblock 收录于《中文工作笔记》v5.2；文集 DOI：\url{https://doi.org/10.5281/zenodo.22739417}。
\newblock 本地源文稿：\path{docs/project_docs/notebook_tex_v5.2/20260313_gframework_discrete_repair_tower_quantum_geometry_interface_formal_system.tex}，第1--3188行。
\bibitem{S22}
高政（GaoZheng）. 基于扩展签发协议的主丛联络--路径积分--同伦修复塔 与结构表示定理. 2026-03-13.
\newblock 收录于《中文工作笔记》v5.2；文集 DOI：\url{https://doi.org/10.5281/zenodo.22739417}。
\newblock 本地源文稿：\path{docs/project_docs/notebook_tex_v5.2/20260313_gframework_extended_certification_tower.tex}，第1--1038行。
\bibitem{S23}
高政（GaoZheng）. 纤维--截面--孔洞--障碍量--势函数--条件化收敛 （独立说明稿：与层 $l_n$ 同伦塔链的对齐与收敛模板，形式逻辑版）. 2026-03-13.
\newblock 收录于《中文工作笔记》v5.2；文集 DOI：\url{https://doi.org/10.5281/zenodo.22739417}。
\newblock 本地源文稿：\path{docs/project_docs/notebook_tex_v5.2/20260313_gframework_fiber_section_hole_obstruction_potential_conditional_convergence_formal_system.tex}，第1--1722行。
\bibitem{S24}
高政（GaoZheng）. 离散拓扑基底上的 Jacobi--Bianchi 障碍修复、 可行路径空间、离散路径积分与终点正规形 （工作笔记：形式系统版）. 2026-03-15.
\newblock 收录于《中文工作笔记》v5.2；文集 DOI：\url{https://doi.org/10.5281/zenodo.22739417}。
\newblock 本地源文稿：\path{docs/project_docs/notebook_tex_v5.2/20260315_gframework_discrete_jacobi_bianchi_repair_tower_path_integral_terminal_normal_form_formal_system.tex}，第1--7415行。
\bibitem{S25}
高政（GaoZheng）. V\(_2\)/V\(_3\) 从 GRL-SAC 中心到开放系统博弈修复引擎族的统一重构 （工作笔记：形式系统版）. 2026-03-18.
\newblock 收录于《中文工作笔记》v5.2；文集 DOI：\url{https://doi.org/10.5281/zenodo.22739417}。
\newblock 本地源文稿：\path{docs/project_docs/notebook_tex_v5.2/20260318_gframework_v2_v3_open_system_game_statistical_repair_family_formal_system.tex}，第1--4416行。
\bibitem{S26}
高政（GaoZheng）. L1 离散、L2 连续与 $L_{\ge 3}$ 同伦修复塔 的纤维丛正规型与孔洞修复 （工作笔记：形式系统版）. 2026-03-27.
\newblock 收录于《中文工作笔记》v5.2；文集 DOI：\url{https://doi.org/10.5281/zenodo.22739417}。
\newblock 本地源文稿：\path{docs/project_docs/notebook_tex_v5.2/20260327_gframework_l1_l2_l3_homotopy_repair_bundle_formal_system.tex}，第1--5110行。
\bibitem{S27}
高政（GaoZheng）. 开放异构系统的宏观拓扑--微观跃迁双层驱动、 广义博弈与总复形同调 （工作笔记：形式系统版）. 2026-04-03.
\newblock 收录于《中文工作笔记》v5.2；文集 DOI：\url{https://doi.org/10.5281/zenodo.22739417}。
\newblock 本地源文稿：\path{docs/project_docs/notebook_tex_v5.2/20260403_gframework_open_heterogeneous_system_macro_micro_total_homology_formal_system.tex}，第1--7655行。
\bibitem{S28}
高政（GaoZheng）. 钱学森弹道、钱学森障碍与观测--推断--障碍修复： 能力元组、边缘算子与全局截面障碍 （工作笔记：形式系统版）. 2026-04-13.
\newblock 收录于《中文工作笔记》v5.2；文集 DOI：\url{https://doi.org/10.5281/zenodo.22739417}。
\newblock 本地源文稿：\path{docs/project_docs/notebook_tex_v5.2/20260413_gframework_observation_inference_obstruction_edge_operator_global_section_formal_system.tex}，第1--10770行。
\bibitem{S29}
高政（GaoZheng）. G框架中性变态射力迫属性、\texorpdfstring{$L_\infty$}{L-infinity}无故障修复与三层证明强度 （工作笔记：形式系统版）. 2026-04-27.
\newblock 收录于《中文工作笔记》v5.2；文集 DOI：\url{https://doi.org/10.5281/zenodo.22739417}。
\newblock 本地源文稿：\path{docs/project_docs/notebook_tex_v5.2/20260427_gframework_forcing_property_morphism_linfty_failure_free_repair_formal_system.tex}，第1--3629行。
\bibitem{S30}
高政（GaoZheng）. 绝对逻辑锚点、粗暴超级对齐与 LLM 日常类幻觉的 语义孔洞及边缘算子修复 （工作笔记：形式系统版）. 2026-04-28.
\newblock 收录于《中文工作笔记》v5.2；文集 DOI：\url{https://doi.org/10.5281/zenodo.22739417}。
\newblock 本地源文稿：\path{docs/project_docs/notebook_tex_v5.2/20260428_gframework_truth_anchor_superalignment_hallucination_obstruction_formal_system.tex}，第1--4158行。
\bibitem{S31}
高政（GaoZheng）. 拒绝签发、概念降级与语义附加：AI 协作中的修复算子注入机制 （工作笔记：形式系统版）. 2026-05-01.
\newblock 收录于《中文工作笔记》v5.2；文集 DOI：\url{https://doi.org/10.5281/zenodo.22739417}。
\newblock 本地源文稿：\path{docs/project_docs/notebook_tex_v5.2/20260501_gframework_refusal_issuance_concept_downgrade_semantic_adjunction_formal_system.tex}，第1--764行。
\bibitem{S32}
高政（GaoZheng）. 统计力学参数格点、同伦扭曲命中与 PDE 策略簇： 从时域求解到频域簇命中的近似量子软件层算法 （工作笔记：形式系统版）. 2026-05-03.
\newblock 收录于《中文工作笔记》v5.2；文集 DOI：\url{https://doi.org/10.5281/zenodo.22739417}。
\newblock 本地源文稿：\path{docs/project_docs/notebook_tex_v5.2/20260503_gframework_statistical_lattice_homotopy_tunnel_pde_strategy_cluster_formal_system.tex}，第1--2868行。
\bibitem{S33}
高政（GaoZheng）. 基于高阶正交确定性收敛的统计力学同伦命中：PDE 求解范式替代 （工作笔记：形式系统版）. 2026-05-13.
\newblock 收录于《中文工作笔记》v5.2；文集 DOI：\url{https://doi.org/10.5281/zenodo.22739417}。
\newblock 本地源文稿：\path{docs/project_docs/notebook_tex_v5.2/20260513_gframework_homotopy_hitting_pde_paradigm_replacement_formal_system.tex}，第1--12938行。
\bibitem{S34}
高政（GaoZheng）. LHTS：延迟信用分配的格点同伦扭曲评分与命中证书体系 （工作笔记：形式系统版）. 2026-05-14.
\newblock 收录于《中文工作笔记》v5.2；文集 DOI：\url{https://doi.org/10.5281/zenodo.22739417}。
\newblock 本地源文稿：\path{docs/project_docs/notebook_tex_v5.2/20260514_gframework_lhts_delayed_credit_lattice_homotopy_twist_scoring_formal_system.tex}，第1--1204行。
\bibitem{S35}
高政（GaoZheng）. TPH-LHTS 轨迹函数参数同伦选择： 概率路径同伦类、广义分形轨迹族与 D 结构递归决策路径 （工作笔记：形式系统版）. 2026-05-15.
\newblock 收录于《中文工作笔记》v5.2；文集 DOI：\url{https://doi.org/10.5281/zenodo.22739417}。
\newblock 本地源文稿：\path{docs/project_docs/notebook_tex_v5.2/20260515_gframework_tph_lhts_trajectory_parameter_homotopy_probability_fractal_d_structure_formal_system.tex}，第1--8054行。
\bibitem{S36}
高政（GaoZheng）. 从 PDE 残差下降到 L-测度微观隧穿： 表达替代、求解替代与解释等价 的形式逻辑证书 （工作笔记：形式系统版）. 2026-05-18.
\newblock 收录于《中文工作笔记》v5.2；文集 DOI：\url{https://doi.org/10.5281/zenodo.22739417}。
\newblock 本地源文稿：\path{docs/project_docs/notebook_tex_v5.2/20260518_gframework_pde_lmeasure_tunnel_equiv_formal_system.tex}，第1--6371行。
\bibitem{S37}
高政（GaoZheng）. G 框架中的数学、物理、计算与认知基础障碍谱系： 统一法空间、总同调语义场与证书化化解 （工作笔记：形式系统版）. 2026-07-08.
\newblock 收录于《中文工作笔记》v5.2；文集 DOI：\url{https://doi.org/10.5281/zenodo.22739417}。
\newblock 本地源文稿：\path{docs/project_docs/notebook_tex_v5.2/20260708_gframework_foundational_obstruction_spectrum_formal_system.tex}，第1--10126行。
\bibitem{S38}
高政（GaoZheng）. 率与商：G 框架语义微积分与传统微积分之间的 同伦映射形式系统. 2026-09-12.
\newblock 收录于《中文工作笔记》v5.2；文集 DOI：\url{https://doi.org/10.5281/zenodo.22739417}。
\newblock 本地源文稿：\path{docs/project_docs/notebook_tex_v5.2/20260912_gframework_rate_quotient_semantic_calculus_traditional_calculus_homotopy_mapping_formal_system.tex}，第1--3934行。
\end{thebibliography}
\clearpage
\end{CJK*}
\end{document}
